%!TeX root=csp-proof.tex
\section{Proofs of the strong-subalgebra theory}\label{sec:strongproofs}

This section proves the statements of \S\ref{sec:properties}. Throughout,
Convention~\ref{conv:standing} is in force --- every algebra is finite,
idempotent, and has a WNU term operation --- and \S\ref{sec:legislation} fixes
the readings.

\subsection{What is imported}\label{sec:imports}

The development is self-contained except for the following, which are proved in
the literature and used here as black boxes. Two are substantial; the rest are
routine. \S\ref{sec:imports-ledger} lists them again with their exact
hypotheses.

\begin{longtheorem}[Absorption Theorem]\label{thm:brady-absorption}
Let $R\sdle\mathbf A\times\mathbf B$ be a subdirect relation between finite
idempotent Taylor algebras. If $R$ is linked then one of the following holds:
$R=A\times B$; or $\mathbf A$ has a proper nonempty binary absorbing or
centrally absorbing subalgebra; or $\mathbf B$ has a proper nonempty subalgebra
that is both binary absorbing and centrally absorbing.
\end{longtheorem}

\begin{corollary}[{a corollary of Theorem~\ref{thm:brady-absorption}}]
\label{lem:linked-implies-abs}
Let $R\lneq_{sd}\mathbf A\times\mathbf B$ be a \emph{proper} subdirect
relation between finite idempotent Taylor algebras. If $R$ is linked then
$\mathbf A$ or $\mathbf B$ has a proper nonempty binary absorbing or central
subuniverse.
\end{corollary}

\begin{proof}
Properness excludes the first alternative of
Theorem~\ref{thm:brady-absorption}. The second is the conclusion on
$\mathbf A$; the third is stronger than the conclusion on $\mathbf B$.
\end{proof}

\begin{hazard}[Both words in ``proper nonempty'' are load-bearing]
\label{haz:linked-abs-proper}
Without them Corollary~\ref{lem:linked-implies-abs} is vacuous, since every
algebra is a binary absorbing subuniverse of itself and $\varnothing$ is
vacuously one (\S\ref{sec:legislation}, item \ref{leg:empty}); every use below
is the proper nonempty one. Theorem~\ref{thm:brady-absorption} is the single
heaviest import of this section, and it is stated here in the form its source
gives rather than in the weaker symmetric form actually used, so that the
strengthening is visible.
\end{hazard}

\begin{lemma}[{\cite[Lem.~2.9]{DecidingAbsorption}}, {\cite[Lem.~6.1, Thm.~6.9]{ZhukStrong}}]\label{lem:pp-propagation}
Let $R\le A^{n}$ be defined by a pp-formula $\Phi$ in which a relation $S$
occurs, and let $R'$ be defined by the formula obtained from $\Phi$ by
replacing \emph{every} occurrence of $S$ by some $S'\subt{T}S$, where
$T\in\{\TBA,\TC\}$. Then $R'\dsubte{T}R$.
\end{lemma}

\begin{hazard}[The conclusion is dotted]\label{haz:pp-dotted}
Replacing a relation by a proper strong subrelation can make the pp-defined
result \emph{empty}, so the undotted $R'\subte{T}R$ --- which under
Convention~\ref{conv:empty} entails $R'\neq\varnothing$ --- is not what the
lemma gives.

Every use below has been checked against this. Lemma~\ref{lem:ba-center-intersect}
states its conclusion dotted and passes it on. Lemmas~\ref{lem:absorb-linked-sub},
\ref{lem:block-good-bridge} and \ref{lem:pc-inductive-step} exhibit a point of
the smaller relation and so upgrade to the undotted form.
Lemma~\ref{lem:bacon-left} needs both nonemptiness \emph{and} properness, and
gets each from a different clause of Lemma~\ref{lem:barto-kazda}
(Caveat~\ref{haz:bacon-left}). The remaining use is inside
Lemma~\ref{lem:center-pushed-in}, which is an outline
(\S\ref{sec:status}).
\end{hazard}

\begin{lemma}[{\cite[Prop.~2.14]{DecidingAbsorption}}, {\cite[Lem.~3.2]{ZhukStrong}}]\label{lem:barto-kazda}
For $B\le\mathbf A$ and $n\ge2$: $B$ absorbs $\mathbf A$ with an operation of
arity $n$ if and only if there is no $S\le A^{n}$ with $S\cap B^{n}=\varnothing$
and $S\cap(B^{i-1}\times A\times B^{n-i})\neq\varnothing$ for every $i$.
\end{lemma}

\begin{lemma}[{\cite[Lem.~6.8]{ZhukStrong}}]\label{lem:factor-type}
If $B\subte{T}\mathbf A$ with $T\in\{\TBA,\TC,\TS\}$ and $\sigma$ is a
congruence on $\mathbf A$, then $B/\sigma\subte{T}\mathbf A/\sigma$.
\end{lemma}

\begin{hazard}[Only the $\TC$ case is imported]\label{haz:factor-type}
For $T=\TBA$ the statement is one line: if $t(B,A)\subseteq B$ and
$t(A,B)\subseteq B$ then the same $t$ works on images. The case $T=\TS$ follows
from the other two by Definition~\ref{def:ba-center-free}. Only $\TC$ is
genuinely imported.
\end{hazard}

\begin{lemma}[{\cite[Lem.~6.24]{ZhukStrong}}]\label{lem:power-to-base}
If $R$ is a nontrivial strong subuniverse of
$\mathbf A_{1}\times\dots\times\mathbf A_{n}$ of a type $T\neq\TPC$, then some
$\mathbf A_{i}$ has a nontrivial subuniverse of type $T$. In particular
$B\subt{T}\mathbf A^{n}$ with $T\in\{\TBA,\TC,\TS\}$ yields some
$C\subt{T}\mathbf A$.
\end{lemma}

\begin{hazard}[One run of the proof serves all three types]\label{haz:power-to-base}
The cited lemma is stated for every type $T\neq\TPC$, and its proof is
\emph{type-uniform}: it takes $\proj_{1}(R)$ if that is proper and otherwise a
fibre $R'=\{(a_{2},\dots,a_{n}):(a_{1},\dots,a_{n})\in R\}$, and neither
construction mentions the type --- the type enters only through the facts that
projections and fibres preserve $\TBA$ and preserve $\TC$. So a single run
produces one witness carrying \emph{every} strong type that $R$ has; in
particular a subuniverse that is simultaneously binary absorbing and central
yields one that is too, which is the case $T=\TS$.
\end{hazard}

\begin{longtheorem}[{\cite[Thm.~6.15]{ZhukStrong}}]\label{thm:central-relation-source}
Let $R\sdle\mathbf A\times\mathbf B$ and
$C=\{c\in A:\{c\}\times B\subseteq R\}$. Then $C$ is a central subuniverse of
$\mathbf A$, or $\mathbf B$ has a nontrivial binary absorbing subuniverse, or
$\mathbf B$ has a nontrivial \emph{projective} subuniverse.
\end{longtheorem}

\begin{hazard}[The heaviest import after Theorem~\ref{thm:brady-absorption}]
\label{haz:615-todo}
Theorem~\ref{thm:central-relation-source} is taken on faith here. Its proof in
\cite{ZhukStrong} is two pages and uses only the Barto--Kazda criterion
(Lemma~\ref{lem:barto-kazda}) and the theory of projective subuniverses; none
of that theory is needed anywhere else below, which is why absorbing it has
been deferred. The same applies to \cite[Lem.~3.4]{ZhukStrong}, used once, in
the proof of Lemma~\ref{lem:central-relation}.
\end{hazard}

\begin{lemma}[{\cite{HobbyMcKenzie}}]\label{lem:hobby-mckenzie}
An algebra is Abelian if and only if some congruence of its square has the
diagonal as a block; and a finite algebra with a WNU term operation is Abelian
if and only if it is affine.
\end{lemma}

\begin{lemma}\label{lem:bridge-abelian}
Let $|A|>1$. Then $\mathbf A$ is Abelian if and only if there is a bridge
$\delta$ from $0_{\mathbf A}$ to $0_{\mathbf A}$ with
$\bcol\delta=\proj_{1,2}(\delta)=\proj_{3,4}(\delta)=A^{2}$.
\end{lemma}

\begin{proof}
If $\mathbf A$ is Abelian, a congruence of $\mathbf A^{2}$ having the diagonal
as a block is such a bridge; clause (c) needs $|A|>1$, since for $|A|=1$ no
relation properly contains $0_{\mathbf A}=A^{2}$.

Conversely, regard $\delta$ as a binary relation $\hat\delta$ on $A^{2}$,
$\hat\delta\bigl((x_{1},x_{2}),(x_{3},x_{4})\bigr)\iff\delta(x_{1},x_{2},
x_{3},x_{4})$, and put $\rho=\hat\delta\circ\hat\delta^{-1}$, a reflexive
symmetric compatible relation on $\mathbf A^{2}$ --- reflexive because
$\proj_{1,2}(\delta)=A^{2}$ makes $\hat\delta$ left-total. Its powers increase,
so $\theta=\rho^{N}$ for $N\ge|A|^{2}$ is reflexive, symmetric and transitive,
hence a congruence of $\mathbf A^{2}$. Clause (d) of
Definition~\ref{def:bridge} is preserved by each composition, so every pair of
$\theta$ relates a diagonal element only to diagonal elements: the diagonal is
a $\theta$-block. Now apply Lemma~\ref{lem:hobby-mckenzie}.
\end{proof}

\subsection{Absorption preliminaries}\label{sec:abs-prelim}

\begin{lemma}\label{lem:ba-center-intersect}
If $B\subte{T}\mathbf A$ with $T\in\{\TBA,\TC\}$ and $C\le\mathbf A$, then
$B\cap C\dsubte{T}C$.
\end{lemma}

\begin{proof}
$C$ is defined by the pp-formula $A(x)\wedge C(x)$; replacing the occurrence of
the full unary relation $A$ by $B$ yields $B\cap C$, so
Lemma~\ref{lem:pp-propagation} applies.
\end{proof}

\begin{lemma}[No absorbing diagonal]\label{lem:no-absorbing-diagonal}
Let $\mathbf C$ be a finite idempotent algebra with $|C|\ge 2$. Then
$0_{\mathbf C}$ is not an absorbing subuniverse of $\mathbf C^{2}$, of any
arity; in particular $0_{\mathbf C}\not\subte{\TBA}\mathbf C^{2}$ and
$0_{\mathbf C}\not\subte{\TC}\mathbf C^{2}$.
\end{lemma}

\begin{proof}
Suppose $t$ is $k$-ary and absorbs at every position. Fix $i$, take
$a_{j}\in C$ for $j\neq i$ and $b,c\in C$, and apply $t$ coordinatewise to the
tuples $(a_{j},a_{j})\in 0_{\mathbf C}$ for $j\neq i$ and $(b,c)\in C^{2}$ at
position $i$. The result
$\bigl(t(\bar a;b),\,t(\bar a;c)\bigr)$ lies in $0_{\mathbf C}$, so
$t(\bar a;b)=t(\bar a;c)$. As $b,c$ were arbitrary, $t$ does not depend on its
$i$-th argument; and this holds for every $i$, so $t$ is constant. Idempotence
then forces $|C|=1$. A central subuniverse is absorbing by
Definition~\ref{def:central}, so the central case is included.
\end{proof}

\begin{lemma}[Absorption preserves linkedness, subrelation form]\label{lem:absorb-linked-sub}
Let $W\sdle\mathbf P\times\mathbf A$ be linked and let
$\varnothing\neq W'\subte{T}W$ with $T\in\{\TBA,\TC\}$. Put
$P'=\proj_{1}(W')$ and $A'=\proj_{2}(W')$. Then $W'\sdle\mathbf P'\times
\mathbf A'$ is linked.
\end{lemma}

\begin{proof}
Let $\Phi_{k}(x,y)$ be the pp-formula in one binary relation symbol obtained by
iterating $x\sim y\iff\exists z\,(W(x,z)\wedge W(y,z))$ $k$ times. Because $W$
is subdirect, $W\circ W^{-1}$ is reflexive and symmetric on $P$, so for
$k\ge|P|$ the formula $\Phi_{k}$ interpreted in $W$ defines the transitive
closure of $W\circ W^{-1}$, which is $P^{2}$ since $W$ is linked. Interpreted
in $W'$ and for the same $k\ge|P|\ge|P'|$ it defines
$\lambda:=\LeftLinked(W')$. By Lemma~\ref{lem:pp-propagation},
replacing every occurrence of $W$ by $W'$ gives $\lambda\dsubte{T}P^{2}$, and
$\lambda\neq\varnothing$ because $W'\neq\varnothing$ makes $P'$ nonempty and
$\lambda$ reflexive on it, so $\lambda\subte{T}P^{2}$. Since
$\lambda\subseteq (P')^{2}\le P^{2}$, Lemma~\ref{lem:ba-center-intersect} gives
$\lambda\subte{T}(P')^{2}$.

Now $\lambda$ is a congruence on $\mathbf P'$: reflexive because $W'$ is
subdirect over $P'$, symmetric by construction, transitive because the powers
stabilise. So $\lambda\times\lambda$ is a congruence on the algebra $(P')^{2}$,
and the image of $\lambda$ under $(P')^{2}\twoheadrightarrow(P'/\lambda)^{2}$ is
exactly $0_{\mathbf P'/\lambda}$; by Lemma~\ref{lem:factor-type},
$0_{\mathbf P'/\lambda}\subte{T}(\mathbf P'/\lambda)^{2}$. Lemma
\ref{lem:no-absorbing-diagonal} forces $|P'/\lambda|=1$, i.e.\
$\lambda=(P')^{2}$, which says $W'$ is linked.
\end{proof}

\begin{hazard}[The box form is not enough]\label{haz:box-vs-sub}
The familiar form of this principle, \cite[Prop.~2.15(i)]{BartoKozik},
restricts by a \emph{box}: if $W\sdle\mathbf A_{1}\times\mathbf A_{2}$ is
linked, $B_{i}$ absorbs $\mathbf A_{i}$, and $W\cap(B_{1}\times B_{2})$ is
subdirect in $B_{1}\times B_{2}$, then $W\cap(B_{1}\times B_{2})$ is linked.
That form cannot reach the restrictions used in
Lemma~\ref{lem:pc-inductive-step}, and the two really do differ: for the
meet-semilattice on $\{0,1,2\}$ with $B=\{0,1\}$, the subalgebra
$K=\{(0,0),(1,2)\}$ of $A^{2}$ meets $B^{2}$ and has
$\proj_{1}(K)\cap B=\{0,1\}$ while $\proj_{1}(K\cap B^{2})=\{0\}$. The
subrelation form also discharges the box form's subdirectness hypothesis for
free.
\end{hazard}

\begin{lemma}[A WNU forbids essentially unary quotients]\label{lem:no-unary-quotient}
Let $\mathbf U$ be a finite algebra with a WNU term operation and $|U|\ge2$.
Then some term operation of $\mathbf U$ depends on more than one variable.
Consequently, if $\mathbf A$ has a WNU term operation then no
$\mathbf U\in\mathrm{HS}(\mathbf A)$ with $|U|\ge2$ is essentially unary.
\end{lemma}

\begin{proof}
Suppose every term operation of $\mathbf U$ has at most one non-dummy variable,
and let $w$ be a WNU term operation of arity $k\ge2$; it is idempotent. If $w$
has no non-dummy coordinate it is constant, and idempotence gives $|U|=1$.
Otherwise $w(\bar x)=g(x_{i})$ for a unary $g$, and $g(x)=w(x,\dots,x)=x$, so
$w$ is the $i$-th projection. The WNU identities equate the $k$ terms in which
a single $y$ sits at each position against a background of $x$'s; two of them
are $w$ with $y$ at position $i$, which evaluates to $y$, and $w$ with $y$ at
some position $\neq i$, which evaluates to $x$. Hence $x=y$ for all $x,y$ and
$|U|=1$. The consequence follows because idempotence and the WNU identities are
identities, hence hold in every member of $\mathrm{HS}(\mathbf A)$.
\end{proof}

\begin{longlemma}[Central relations]\label{lem:central-relation}
Let $\mathbf A,\mathbf B$ satisfy Convention~\ref{conv:standing},
let $R\sdle\mathbf A\times\mathbf B$, and put
$C=\{c\in A:\{c\}\times B\subseteq R\}$. Then
\begin{enumerate}\alphlabels\tightlist
\item $C=\varnothing$, or
\item $C$ is a central subuniverse of $\mathbf A$, or
\item $\mathbf B$ has a proper nonempty binary absorbing subuniverse.
\end{enumerate}
\end{longlemma}

\begin{proof}
Apply Theorem~\ref{thm:central-relation-source}. Its first alternative gives
(a) or (b) according as $C$ is empty or not, its second gives (c). In its third,
$\mathbf B$ has a nontrivial projective subuniverse $P$; by
\cite[Lem.~3.4]{ZhukStrong}, either $P$ is a binary absorbing subuniverse of
$\mathbf B$, which is (c) since $P$ is nontrivial, or there is an essentially
unary $\mathbf U\in\mathrm{HS}(\mathbf B)$ with $|U|\ge2$, which
Lemma~\ref{lem:no-unary-quotient} forbids.
\end{proof}

\begin{hazard}[Alternative (a) is not optional]\label{haz:central-relation}
Under Convention~\ref{conv:empty} the two-alternative form (b)--(c) is false:
$\mathbf A=\mathbf B=\mathbf Z_{p}$ with $R=\{(x,x+1)\}$ is subdirect, has
$C=\varnothing$, and $\mathbf Z_{p}$ has no nontrivial binary absorbing
subuniverse at all. It is true in \cite{ZhukStrong} only because $\varnothing$
counts as central there. Every call site below discharges (a) by exhibiting a
point of $C$ first, and two further obligations go with each contrapositive
use: the centre produced must be shown \emph{proper}, and a strong subuniverse
of a \emph{power} must be moved down by Lemma~\ref{lem:power-to-base}.
\end{hazard}

\begin{lemma}\label{lem:bacon-left}
Let $R\sdle\mathbf A\times\mathbf B$ with $\mathbf A$ BA and centre free,
$\LeftLinked(R)=A^{2}$, and
$C=\{c\in B:A\times\{c\}\subseteq R\}$. Then $C\neq\varnothing$ and
$C\subte{\TBA,\TC}\mathbf B$.
\end{lemma}

\begin{proof}
If $|A|=1$ then subdirectness gives $A\times\{b\}\subseteq R$ for every
$b\in B$, so $C=B$ and the conclusion holds with equality. Assume $|A|\ge2$.

\emph{Step 1: $C\neq\varnothing$.} For $k\ge1$ put
$W_{k}=\{(a_{1},\dots,a_{k}):\exists b\,\forall i\,R(a_{i},b)\}$, a
subuniverse of $\mathbf A^{k}$, with $W_{1}=\proj_{1}(R)=A$. If
$W_{|A|}=A^{|A|}$ then a witness $b$ for the tuple listing all of $A$ has
$A\times\{b\}\subseteq R$, so $b\in C$ and we are done. Otherwise choose
$k\ge2$ minimal with $W_{k}\neq A^{k}$, so $W_{k-1}=A^{k-1}$.

View $W_{k}$ as a binary relation in $\mathbf A\times\mathbf A^{k-1}$. It is
\emph{subdirect}: for $a\in A$, subdirectness of $R$ gives $b$ with $R(a,b)$
and hence $(a,a,\dots,a)\in W_{k}$, so the left projection is $A$; and for
$\bar a\in A^{k-1}$, $W_{k-1}=A^{k-1}$ gives $b$ with $R(a_{i},b)$ for all
$i$, while $\proj_{2}(R)=B$ gives some $a_{0}$ with $R(a_{0},b)$, so
$(a_{0},\bar a)\in W_{k}$ and the right projection is $A^{k-1}$.

It is \emph{linked}. If $a,c\in A$ share an $R$-neighbour $b$ --- that is,
$R(a,b)$ and $R(c,b)$ --- then $b$ also witnesses
$(a,a,\dots,a),(c,a,\dots,a)\in W_{k}$, so $(a,c)\in\LeftLinked(W_{k})$ by
Definition~\ref{def:linked}. Since $\LeftLinked(R)$ is the least equivalence
closed under that step and $\LeftLinked(W_{k})$ is an equivalence containing
each such pair, $\LeftLinked(R)\subseteq\LeftLinked(W_{k})$; and
$\LeftLinked(R)=A^{2}$ by hypothesis, so $\LeftLinked(W_{k})=A^{2}$.

So $W_{k}$ is a proper linked subdirect relation, and
Corollary~\ref{lem:linked-implies-abs} gives a proper nonempty binary absorbing
or central subuniverse of $\mathbf A$ or of $\mathbf A^{k-1}$;
Lemma~\ref{lem:power-to-base} moves the second case down to $\mathbf A$. Either
way this contradicts the hypothesis on $\mathbf A$. Hence
$W_{|A|}=A^{|A|}$ after all, and $C\neq\varnothing$.

\emph{Step 2: $C$ is central in $\mathbf B$.} Apply
Lemma~\ref{lem:central-relation} to $R^{-1}\sdle\mathbf B\times\mathbf A$,
whose associated set is exactly $C$. Alternative (a) is excluded by Step 1 and
alternative (c) --- a proper nonempty binary absorbing subuniverse of
$\mathbf A$ --- by hypothesis, so (b) holds.

\emph{Step 3: $C$ is binary absorbing in $\mathbf B$.} Suppose not. By
Lemma~\ref{lem:barto-kazda} at $n=2$ there is $S\le B\times B$ with
\[
  S\cap(C\times C)=\varnothing,\qquad
  S\cap(C\times B)\neq\varnothing,\qquad
  S\cap(B\times C)\neq\varnothing .
\]
Write $U_{X}=\{b\in B:\exists c\,(S(b,c)\wedge X(c))\}$ for a unary $X$ on
$B$, and consider the pp-formula
\[
  \Phi_{X}(a_{1},\dots,a_{|A|})=\exists b\;\Bigl(U_{X}(b)\wedge
    \bigwedge_{i}R(a_{i},b)\Bigr),
\]
in which the unary relation $X$ occurs once. Put $W=\Phi_{C}$ and
$W^{\circ}=\Phi_{B}$. Since $C\subt{\TC}\mathbf B$ by Step 2 ---
\emph{proper}, because $C=B$ would put $S\cap(C\times C)=S\neq\varnothing$ ---
Lemma~\ref{lem:pp-propagation} gives $W\dsubte{\TC}W^{\circ}$. Three
observations turn that into what is wanted, and each spends one of the three
displayed conditions.

\emph{$W^{\circ}=A^{|A|}$.} $S\cap(C\times B)\neq\varnothing$ supplies
$b_{0}\in C$ with $b_{0}\in U_{B}$. Since $b_{0}\in C$ we have
$A\times\{b_{0}\}\subseteq R$, so $b_{0}$ witnesses $\Phi_{B}(\bar a)$ for
\emph{every} $\bar a\in A^{|A|}$.

\emph{$W\neq\varnothing$.} $S\cap(B\times C)\neq\varnothing$ supplies
$(b_{1},c_{1})\in S$ with $c_{1}\in C$, so $b_{1}\in U_{C}$; and
$\proj_{2}(R)=B$ supplies $a$ with $R(a,b_{1})$, so
$(a,\dots,a)\in W$.

\emph{$W\neq A^{|A|}$.} If it were, the tuple enumerating $A$ would give
$b\in U_{C}$ with $A\times\{b\}\subseteq R$, i.e.\ $b\in C$; and $b\in U_{C}$
means $S(b,c)$ for some $c\in C$, so $(b,c)\in S\cap(C\times C)=\varnothing$.

So $W\subt{\TC}A^{|A|}$ is proper and nonempty, and
Lemma~\ref{lem:power-to-base} produces a proper nonempty central subuniverse of
$\mathbf A$ --- again a contradiction.
\end{proof}

\begin{hazard}[All three Barto--Kazda conditions are spent]\label{haz:bacon-left}
Step 3 is the only place in this section where the three side conditions of
Lemma~\ref{lem:barto-kazda} are used separately, one apiece: $S$ meeting
$C\times B$ is what makes the \emph{ambient} of the pp-propagation step the
full power $A^{|A|}$ rather than some subrelation of it; $S$ meeting
$B\times C$ is what makes $W$ nonempty; and $S\cap(C\times C)=\varnothing$ is
what makes it proper. Dropping any one leaves a conclusion that
Lemma~\ref{lem:power-to-base} cannot consume, since $\subt{\TC}$ requires a
proper nonempty subuniverse of a product.
\end{hazard}

\begin{lemma}\label{lem:strong-nonempty}
If $C\lll^{A}B\lll A$ and $D\subt{\TBA,\TC}B$ then $C\cap D\neq\varnothing$.
\end{lemma}

\begin{proof}
Suppose not, and choose $C''$ minimal with
$C\lll^{A}C'\subt[A]{T(\sigma)}C''\lll^{A}B$ and $C''\cap D\neq\varnothing$.
By Lemma~\ref{lem:ba-center-intersect} $C''\cap D\subt{\TBA,\TC}C''$, so by
Lemma~\ref{lem:possible-intersections} $T=\TD$; and then
Lemma~\ref{lem:factor-type} gives
$(C''\cap D)/\sigma\subt{\TBA}C''/\sigma$, contradicting the requirement in
Definition~\ref{def:types} that $C''/\sigma$ be BA and centre free.
\end{proof}

\begin{lemma}[No critical central family of size $\ge3$]\label{lem:no-central-critical}
Let $n\ge3$, let $B_{1},\dots,B_{n}\le\mathbf A$, and let
$C_{i}\subt{\TC}B_{i}$ for $i\in[n]$. Then
$\bigcap_{i}C_{i}=\varnothing$ and $B_{j}\cap\bigcap_{i\neq j}C_{i}\neq
\varnothing$ for every $j$ cannot both hold.
\end{lemma}

\begin{proof}
Suppose they do. Put $D=\bigcap_{i}B_{i}$ and $E_{i}=C_{i}\cap D$. Then $D$ is
a nonempty subuniverse, containing $B_{1}\cap C_{2}\cap\dots\cap C_{n}$; each
$E_{i}\subte{\TC}D$ by Lemma~\ref{lem:ba-center-intersect}, and properly, since
$E_{j}=D$ would give $\bigcap_{i}E_{i}=\bigcap_{i\neq j}E_{i}\neq\varnothing$.
Moreover $\bigcap_{i}E_{i}=\varnothing$ and, because $C_{i}\subseteq B_{i}$,
\[
  \textstyle\bigcap_{i\neq j}E_{i}=D\cap\bigcap_{i\neq j}C_{i}
  =B_{j}\cap\bigcap_{i\neq j}C_{i}\neq\varnothing .
\]
Now the diagonal $R=\{(a,\dots,a):a\in D\}\le D^{n}$ is a subuniverse by
idempotence, and the two displays say exactly that $R$ is
$(E_{1},\dots,E_{n})$-essential. By \cite[Lem.~6.11]{ZhukStrong} no such
relation exists for $n\ge3$ when the $E_{i}$ are proper nonempty central
subuniverses.
\end{proof}

\begin{corollary}[Relational form]\label{cor:no-central-critical}
Let $R\sdle\mathbf A_{1}\times\dots\times\mathbf A_{n}$ with $n\ge3$,
$B_{i}\le\mathbf A_{i}$, and $C_{i}\subt{\TC}B_{i}$ for $i\in[n]$. If
$R\cap(C_{1}\times\dots\times B_{j}\times\dots\times C_{n})\neq\varnothing$
for every $j$, then $R\cap(C_{1}\times\dots\times C_{n})\neq\varnothing$.
\end{corollary}

\begin{proof}
Write $f_{i}\colon\mathbf R\to\mathbf A_{i}$ for the projections and
$P_{i}=f_{i}^{-1}(B_{i})$, $Q_{i}=f_{i}^{-1}(C_{i})$, subuniverses of
$\mathbf R$. The formula $R(x_{1},\dots,x_{n})\wedge X(x_{i})$ is positive
primitive in $X$ and defines $P_{i}$ at $X=B_{i}$ and $Q_{i}$ at $X=C_{i}$, so
Lemma~\ref{lem:pp-propagation} gives $Q_{i}\dsubte{\TC}P_{i}$; and $Q_{i}$ is
nonempty, containing $R\cap(C_{1}\times\dots\times B_{j}\times\dots\times
C_{n})$ for any $j\neq i$, which exists since $n\ge3$. So
$Q_{i}\subte{\TC}P_{i}$.

If $Q_{i}=P_{i}$ for some $i$, the conclusion is immediate:
$R\cap(C_{1}\times\dots\times C_{n})
=Q_{i}\cap\bigcap_{j\neq i}Q_{j}
=P_{i}\cap\bigcap_{j\neq i}Q_{j}
=R\cap(C_{1}\times\dots\times B_{i}\times\dots\times C_{n})\neq\varnothing$.
Otherwise every $Q_{i}\subt{\TC}P_{i}$ is proper, and
Lemma~\ref{lem:no-central-critical} applies inside $\mathbf R$ with
$B_{i}:=P_{i}$ and $C_{i}:=Q_{i}$: its two forbidden conditions are exactly
$R\cap(C_{1}\times\dots\times C_{n})=\varnothing$ and the displayed
nonemptiness, read through the $f_{i}$. So the first fails.
\end{proof}

\begin{hazard}[Why this is a separate statement]\label{haz:central-critical-split}
Lemma~\ref{lem:no-central-critical} is the type-$\TC$ half of case (c) of
Theorem~\ref{thm:stable-intersection}, and it is stated on its own because
Lemma~\ref{lem:block-good-bridge} needs \emph{only} this half. Citing the whole
theorem there would make the bridge theory rest on the bridge classification,
which rests back on it:
\[
  \text{\ref{cor:stable-intersection}}\to
  \text{\ref{thm:stable-intersection}}\to
  \text{\ref{lem:bridge-to-pc}}\to\text{\ref{lem:no-cross-bridge}}\to
  \text{\ref{lem:linear-equiv}}\to\text{\ref{lem:nontrivial-bridge}}\to
  \text{\ref{lem:block-good-bridge}}\to\text{\ref{cor:stable-intersection}} .
\]
Every edge of that is a real citation, and the loop is broken only by observing
that the appeal at the last step uses no bridge at all: it is $n=3$ with all
types $\TC$, and the argument for that case runs through
\cite[Lem.~6.11]{ZhukStrong} and nothing else in this section. Splitting the
statement is what makes the observation checkable rather than a matter of
trust.

A second observation, found by the layer check rather than by reading: the
$\lll$ hypotheses these two statements used to carry are consumed by
\emph{nothing} in either proof --- nonemptiness of $\bigcap_{i}B_{i}$ comes
from the displayed conditions, and the essential-relations import needs only
proper nonempty central subuniverses. They are now stated for plain
subalgebras, and the corollary pulls back through
Lemma~\ref{lem:pp-propagation} instead of the $\lll$-pullback of
Lemma~\ref{lem:reverse-hom}, with the equality case of the pullback giving the
conclusion outright. That is what places the entire bridge classification
below the strong-subuniverse theory rather than beside it.
\end{hazard}

\begin{lemma}[{\cite[Lem.~6.25]{ZhukStrong}}]\label{lem:possible-intersections}
If $B\subt{T_{1}}\mathbf A$, $C\subt{T_{2}}\mathbf A$, $B\cap C=\varnothing$
and $T_{1},T_{2}\in\{\TBA,\TC,\TS\}$, then $T_{1}=T_{2}\in\{\TBA,\TC\}$.
\end{lemma}

\begin{lemma}[A BA-and-central subuniverse meets every strong one]\label{lem:ba-central-meets}
If $D\subt{\TBA,\TC}\mathbf A$ and $E\subt{T}\mathbf A$ with
$T\in\{\TBA,\TC,\TS\}$, then $D\cap E\neq\varnothing$.
\end{lemma}

\begin{proof}
Suppose $D\cap E=\varnothing$. Lemma~\ref{lem:possible-intersections} applied
to the pair $(D,E)$ with $T_{1}=\TBA$ gives $T=\TBA$; applied to the same pair
with $T_{1}=\TC$ it gives $T=\TC$. The two conclusions are incompatible.
\end{proof}

\begin{hazard}[Where the conjunction reading of $\subt{\TBA,\TC}$ pays]\label{haz:ba-central-meets}
Lemma~\ref{lem:ba-central-meets} is false with $\subt{\TBA,\TC}$ read as a
disjunction: two disjoint binary absorbing subuniverses are exactly what
Lemma~\ref{lem:possible-intersections} permits. What makes the argument run is
that $D$ carries \emph{both} types at once, so whichever of $\TBA$ and $\TC$
the subuniverse $E$ has, the other is available for $D$ and the two types
differ. This is the reading fixed in Caveat~\ref{haz:S-type}; it is the
mechanism already used inside the proof of Lemma~\ref{lem:strong-nonempty},
and the two appeals in Lemma~\ref{lem:intersection-good}(s) are what make it
worth naming.
\end{hazard}

\subsection{The intersection property}\label{sec:intersection-property}

The goal is Lemma~\ref{lem:intersection-good}: if $B\lll A$, $C\lll A$ meet and
$\delta$ is a dividing congruence for $B$, then $(B\cap C)/\delta$ is a single
point or all of $B/\delta$. Everything in \S\ref{sec:strongproofs} that
produces a bridge goes through it.

\begin{lemma}[The shift manoeuvre]\label{lem:multiset-shift}
Let $k\ge2$ and let $R\le A^{k}$ be totally symmetric and closed under
$(a_{1},\dots,a_{k})\mapsto(a_{1},a_{1},a_{2},\dots,a_{k-1})$. Identify a tuple
of $R$ with its entry multiset, a $k$-element multiset over $A$. If
$\bar a\in R$ has entry multiset $M$, then for any $u,v\in M$ ---
\emph{occupying distinct positions of $\bar a$, though possibly with $u=v$ as
elements of $A$} --- the multiset $M+\{u\}-\{v\}$ is again the entry multiset of
a tuple of $R$. Consequently every $k$-element multiset over the set $S$ of
distinct entries of $\bar a$ is reached, provided it is reached by such steps;
in particular, for any $u\in S$, the constant multiset $k\cdot\{u\}$ is
reached, in at most $k-1$ steps.
\end{lemma}

\begin{proof}
Total symmetry lets us place $u$ first and $v$ last; the shift then duplicates
$u$ and deletes $v$, giving $M+\{u\}-\{v\}$. For the last clause, apply this
$k-1$ times with the same $u$, deleting a different position each time.
\end{proof}

\begin{hazard}[The shift deletes the last entry]\label{haz:shift}
A single application of the shift to $(a,b_{2},\dots,b_{k-1},c)$ destroys $c$,
so $(a,\dots,a,c)\in R$ does not follow from one step; it needs $k-2$ steps
interleaved with permutations, which is what Lemma~\ref{lem:multiset-shift}
packages. The lemma is used twice below and the one-step reading fails both
times.
\end{hazard}

\begin{lemma}\label{lem:tot-sym-full}
Let $k\ge2$, let $\mathbf A$ be BA and centre free, and let $R\le A^{k}$ be
totally symmetric, closed under the shift, with $\proj_{1,2}(R)=A^{2}$. Then
$R=A^{k}$.
\end{lemma}

\begin{proof}
Induct on $k$; for $k=2$ the hypothesis is the conclusion. Let $k>2$.

\emph{Step 1: $(a,\dots,a,c)\in R$ for all $a,c$.} By total symmetry
$\proj_{1,k}(R)=A^{2}$, so there is a tuple of $R$ with first entry $a$ and
last entry $c$. Apply Lemma~\ref{lem:multiset-shift} $k-2$ times with $u=a$,
deleting on each step one position other than the last: the entry multiset
becomes $(k-1)\cdot\{a\}+\{c\}$, which by total symmetry gives
$(a,\dots,a,c)\in R$. Deleting $c$ as well gives $(a,\dots,a)\in R$.

\emph{Step 2: the left projection is full.} $R'=\proj_{1,\dots,k-1}(R)$
inherits total symmetry, has $\proj_{1,2}(R')=A^{2}$, and is shift-closed: from
$\bar a\in R'$ extend to $R$, shift there, and project. By the inductive
hypothesis $R'=A^{k-1}$.

\emph{Step 3.} Only now may $R$ be viewed as a subdirect relation
$R\sdle A^{k-1}\times A$: Step 2 gives the left projection and Step 1 gives the
right. It is linked, since by Step 1 every $c$ is adjacent to the single left
vertex $(a,\dots,a)$. If $R\neq A^{k}$ then
Corollary~\ref{lem:linked-implies-abs} gives a proper nonempty binary
absorbing or central subuniverse on $\mathbf A^{k-1}$ or on $\mathbf A$, and
Lemma~\ref{lem:power-to-base} moves the first case to $\mathbf A$ ---
contradicting the hypothesis either way.
\end{proof}

\begin{lemma}\label{lem:tot-sym-irred}
Let $\sigma$ be a dividing congruence for $B\le A$ and let
$R\le(\mathbf A/\sigma)^{k}$ be reflexive, totally symmetric and shift-closed.
Then $(B/\sigma)^{k}\subseteq R$, or $R$ is the set of constant tuples.
\end{lemma}

\begin{proof}
If $\proj_{1,2}(R)$ is the equality relation then, by total symmetry, all
coordinates of every tuple of $R$ agree, and reflexivity supplies all constant
tuples.

Otherwise $\proj_{1,2}(R)$ is a subuniverse of $(\mathbf A/\sigma)^{2}$
containing $0_{\mathbf A/\sigma}$ properly and stable under it, so by
minimality of the cover (Proposition~\ref{prop:cover}, transported by
Lemma~\ref{lem:irred-quotient}) it contains $\cover{\sigma}/\sigma$, hence
$(B/\sigma)^{2}$ --- the last step by Definition~\ref{def:types}(i). Now apply
Lemma~\ref{lem:tot-sym-full} to $R\cap(B/\sigma)^{k}$ over the algebra
$B/\sigma$, which is BA and centre free by Definition~\ref{def:types}(iii). Its
hypothesis $\proj_{1,2}\bigl(R\cap(B/\sigma)^{k}\bigr)=(B/\sigma)^{2}$ is not
immediate from $\proj_{1,2}(R)\supseteq(B/\sigma)^{2}$, since a witnessing
tuple may leave $B/\sigma$; Lemma~\ref{lem:multiset-shift} repairs it, pulling
$(x_{1},x_{2},y_{3},\dots,y_{k})$ back to $(x_{1},x_{2},x_{1},\dots,x_{1})$.
\end{proof}

\begin{lemma}[Full fibre]\label{lem:full-fibre}
Let $F$ be a finite set with $|F|\le N$ and let $W\subseteq F^{N}\times E$
satisfy $\proj_{1..N}(W)=F^{N}$ and be closed under coordinate substitution: if
$(x_{1},\dots,x_{N},e)\in W$ and $g:[N]\to[N]$ then
$(x_{g(1)},\dots,x_{g(N)},e)\in W$. Then $F^{N}\times\{d\}\subseteq W$ for some
$d\in E$.
\end{lemma}

\begin{proof}
Write $F=\{F_{1},\dots,F_{t}\}$ with $t\le N$ and put
$\bar x=(F_{1},\dots,F_{t},F_{t},\dots,F_{t})\in F^{N}$. By surjectivity pick
$d$ with $(\bar x,d)\in W$. Every $\bar y\in F^{N}$ is $\bar x\circ g$ for some
$g:[N]\to[N]$, since every entry of $\bar y$ occurs among the entries of
$\bar x$; closure gives $(\bar y,d)\in W$.
\end{proof}

\begin{hazard}[Why the arity is $|A|$]\label{haz:full-fibre}
Lemma~\ref{lem:full-fibre} is used three times below, always in the shape
``since $\proj_{1,\dots,|A|}(R')=(B/\delta)^{|A|}$, there is $d$ with
$(B/\delta)^{|A|}\times\{d\}\subseteq R'$''. Surjectivity of a projection does
not by itself produce a full fibre; what makes the step valid is that each such
$R'$ is cut out by a conjunction of unary and binary conditions on its first
$|A|$ entries, hence closed under substituting entries for one another, and
that $|B/\delta|\le|A|$. \textbf{The exponent $|A|$ is load-bearing}: it must be
at least the number of blocks, so that one tuple can enumerate them.
Normalising the arity away breaks all three proofs.
\end{hazard}

\begin{longlemma}\label{lem:self-intersection-pc}
Let $B_{k}\subt[A]{T_{k}(\sigma_{k})}\dots\subt[A]{T_{1}(\sigma_{1})}B_{0}=A$
with $T_{i}\in\{\TBA,\TC,\TS,\TD\}$, let $\delta$ be a congruence on
$\mathbf A$, and let $m\in[k]$ with $T_{m}=\TD$. Then
$\bigl((B_{k}\circ\delta)\cap B_{m-1}\bigr)/\sigma_{m}$ is either
$B_{m-1}/\sigma_{m}$ or $B_{m}/\sigma_{m}$; and in the second case
$\sigma_{m}\supseteq\delta\cap\sigma_{1}\cap\dots\cap\sigma_{m-1}$.
\end{longlemma}

\begin{proof}
The induction is on $k$, and $\delta$, $m$ and the chain are all quantified
\emph{inside} the inductive statement: subsubcase 1B3 below applies the
hypothesis with a different congruence in the role of $\delta$, so $\delta$ must
not be held fixed.

Two conventions are in force throughout. For every $j$ with $T_{j}=\TD$ we have
$B_{j}=B_{j-1}\cap E_{j}$ for a block $E_{j}$ of $\sigma_{j}$, so
$B_{j}^{2}\subseteq\sigma_{j}$, $|B_{j}/\sigma_{j}|=1$, and
$\mathbf B_{j-1}/\sigma_{j}$ is BA and centre free with
$|B_{j-1}/\sigma_{j}|\ge2$ (Definition~\ref{def:types}(iii), and properness of
$B_{j}$ in $B_{j-1}$); in particular $|A|\ge2$. For every other $j$,
$\sigma_{j}=A^{2}$.

\medskip\noindent\textbf{Case 1: $k=m$,} so $T_{k}=\TD$. Put
$\tau=\delta\cap\sigma_{1}\cap\dots\cap\sigma_{k-1}$ and, for
$0\le n\le k$,
\[
  Q_{n}=\bigl\{(a_{1},\dots,a_{|A|})\in B_{n}^{|A|}:
    (a_{i},a_{j})\in\tau\ \text{for all }i,j\bigr\},
  \qquad S_{n}=Q_{n}/\sigma_{k},
\]
the quotient again an image. Each $Q_{n}$ is cut out of $\mathbf A^{|A|}$ by
unary and binary conditions, so each $S_{n}$ is a subuniverse of
$(\mathbf A/\sigma_{k})^{|A|}$ closed under substitution of coordinates and
hence totally symmetric and shift-closed;
$S_{0}$, whose $Q_{0}$ ranges over $B_{0}=A$, is in addition reflexive; and
$S_{n}\supseteq S_{n'}$ for $n\le n'$. Since $\sigma_{k}$ is a dividing
congruence for $B_{k-1}\le A$, Lemma~\ref{lem:tot-sym-irred} applies to
$S_{0}$.

\emph{Subcase 1A: $S_{0}=\{(a/\sigma_{k},\dots,a/\sigma_{k}):a\in A\}$.} For
$(a,b)\in\tau$ the tuple $(a,b,b,\dots,b)$ lies in $Q_{0}$, so its image is
constant and $a/\sigma_{k}=b/\sigma_{k}$: that is
$\sigma_{k}\supseteq\tau$, which is the additional clause. It also settles the
main one. Every $\sigma_{\ell}$ with $\ell<k$ contains $B_{k-1}^{2}$ --- for
$T_{\ell}=\TD$ because $B_{k-1}\subseteq B_{\ell}$ and
$B_{\ell}^{2}\subseteq\sigma_{\ell}$, otherwise because
$\sigma_{\ell}=A^{2}$ --- so
\[
  \delta\cap B_{k-1}^{2}\;\subseteq\;\tau\cap B_{k-1}^{2}\;\subseteq\;\sigma_{k}.
\]
Hence if $a\in B_{k-1}$ and $(a,b)\in\delta$ with $b\in B_{k}\subseteq B_{k-1}$
then $(a,b)\in\sigma_{k}$, so $a\in E_{k}$ and therefore
$a\in B_{k-1}\cap E_{k}=B_{k}$. Thus
$(B_{k}\circ\delta)\cap B_{k-1}=B_{k}$ and the projection is
$B_{k}/\sigma_{k}=B_{m}/\sigma_{m}$.

\emph{Subcase 1B: $(B_{k-1}/\sigma_{k})^{|A|}\subseteq S_{0}$.} Suppose first
that $(B_{k-1}/\sigma_{k})^{|A|}\subseteq S_{k-1}$. Since $|B_{k-1}/\sigma_{k}|
\le|A|$, one tuple $x\in(B_{k-1}/\sigma_{k})^{|A|}$ enumerates every class of
$B_{k-1}/\sigma_{k}$; a witness for $x$ in $Q_{k-1}$ is a family
$a_{1},\dots,a_{|A|}\in B_{k-1}$, pairwise $\tau$-related and so in particular
pairwise $\delta$-related, whose $\sigma_{k}$-classes exhaust
$B_{k-1}/\sigma_{k}$. Let $E$ be the $\delta$-class containing all of them, so
that
\[
  (E\cap B_{k-1})/\sigma_{k}=B_{k-1}/\sigma_{k}.
\]
One of those classes is the class of $B_{k}$, so some $a_{i}$ lies in
$B_{k-1}\cap E_{k}=B_{k}$; hence $E$ meets $B_{k}$, so
$E\subseteq B_{k}\circ\delta$, and
$((B_{k}\circ\delta)\cap B_{k-1})/\sigma_{k}\supseteq(E\cap B_{k-1})/\sigma_{k}
=B_{k-1}/\sigma_{k}$. That is the first alternative, and this branch is done.

So assume $(B_{k-1}/\sigma_{k})^{|A|}\not\subseteq S_{k-1}$ and choose $n$
minimal with $(B_{k-1}/\sigma_{k})^{|A|}\not\subseteq S_{n}$; then
$1\le n\le k-1$ and $(B_{k-1}/\sigma_{k})^{|A|}\subseteq S_{n-1}$. For every
$b\in B_{k-1}\subseteq B_{n}$ the constant tuple $(b,\dots,b)$ lies in
$Q_{n}$, so
\begin{equation}\label{eq:sip-nonempty}
  \varnothing\neq S_{n}\cap(B_{k-1}/\sigma_{k})^{|A|}
  \subsetneq(B_{k-1}/\sigma_{k})^{|A|} .
\end{equation}
We show each of the three types of $T_{n}$ impossible, so that this branch does
not occur.

\emph{Subsubcase 1B1: $T_{n}\in\{\TBA,\TC\}$.} The formula
$\bigwedge_{i,j}\tau(x_{i},x_{j})\wedge\bigwedge_{i}X(x_{i})$ defines
$Q_{n-1}$ at $X=B_{n-1}$ and $Q_{n}$ at $X=B_{n}$, so
Lemma~\ref{lem:pp-propagation} and \eqref{eq:sip-nonempty} give
$Q_{n}\subte{T_{n}}Q_{n-1}$; Lemma~\ref{lem:factor-type} pushes this through
$\sigma_{k}$ and Lemma~\ref{lem:ba-center-intersect} restricts it along
$(B_{k-1}/\sigma_{k})^{|A|}\le\mathbf S_{n-1}$, giving by
\eqref{eq:sip-nonempty} a proper nonempty
$\subt{T_{n}}(B_{k-1}/\sigma_{k})^{|A|}$. Lemma~\ref{lem:power-to-base} moves
it to $B_{k-1}/\sigma_{k}$, which is BA and centre free --- a contradiction.

\emph{Subsubcase 1B2: $T_{n}=\TS$.} Choose a nonempty $G\subseteq B_{n}$ with
$G\subt{\TBA,\TC}B_{n-1}$. Lemma~\ref{lem:strong-nonempty}, applied along
$B_{k}\lll^{A}B_{n-1}\lll A$, gives $G\cap B_{k}\neq\varnothing$, and
$B_{k}\subseteq B_{k-1}$, so replacing $B_{n}$ by $G$ leaves the intersection
with $(B_{k-1}/\sigma_{k})^{|A|}$ nonempty; it stays proper because
$G\subseteq B_{n}$. Now 1B1 applies verbatim with $G$ for $B_{n}$ and
contradicts BA-freeness of $B_{k-1}/\sigma_{k}$.

\emph{Subsubcase 1B3: $T_{n}=\TD$.} Let $R'$ be the image in
$(\mathbf A/\sigma_{k})^{|A|}\times\mathbf A/\sigma_{n}$ of
\[
  \bigl\{(a_{1},\dots,a_{|A|},b):b\in B_{n-1},\ \forall i\;
    a_{i}\in B_{n-1}\cap(B_{k-1}\circ\sigma_{k}),\ \forall i,j\;
    (a_{i},a_{j})\in\tau,\ \forall i\;(a_{i},b)\in\sigma_{n}\bigr\} .
\]
Because $n\le k-1$, the congruence $\sigma_{n}$ is one of those intersected in
$\tau$, so $\tau\subseteq\sigma_{n}$. Four facts.

\emph{(i) $\proj_{1,\dots,|A|}(R')=(B_{k-1}/\sigma_{k})^{|A|}$.} A tuple
$x$ on the right lies in $S_{n-1}$ by minimality of $n$, so it has witnesses
$a_{i}\in B_{n-1}$, pairwise $\tau$-related, with $a_{i}/\sigma_{k}=x_{i}$ and
hence $a_{i}\in B_{k-1}\circ\sigma_{k}$; and $\tau\subseteq\sigma_{n}$ makes
$b=a_{1}\in B_{n-1}$ a legitimate last entry.

\emph{(ii) $\proj_{|A|+1}(R')=((B_{k-1}\circ\sigma_{k})\cap B_{n-1})/\sigma_{n}$.}
A witnessing tuple has $(a_{1},b)\in\sigma_{n}$ with
$a_{1}\in B_{n-1}\cap(B_{k-1}\circ\sigma_{k})$; conversely such a $c$ is
witnessed by $(c,\dots,c,c)$.

\emph{(iii) $(B_{k-1}/\sigma_{k})^{|A|}\times B_{n}/\sigma_{n}\not\subseteq R'$.}
Otherwise, fixing $b_{0}\in B_{k-1}\subseteq B_{n}$, every
$x\in(B_{k-1}/\sigma_{k})^{|A|}$ has a witness with $(a_{i},b_{0})\in
\sigma_{n}$; then $b_{0}\in B_{n}\subseteq E_{n}$ puts each $a_{i}$ in
$B_{n-1}\cap E_{n}=B_{n}$, so $x\in S_{n}$, contradicting the choice of $n$.

\emph{(iv) $\proj_{|A|+1}(R')=B_{n-1}/\sigma_{n}$.} The inductive hypothesis at
$k-1$ --- applied to the chain $B_{k-1}\subt[A]{T_{k-1}(\sigma_{k-1})}\dots
\subt[A]{T_{1}(\sigma_{1})}B_{0}$, the congruence $\sigma_{k}$ in the role of
$\delta$, and the index $n$ --- leaves $B_{n-1}/\sigma_{n}$ or
$B_{n}/\sigma_{n}$ for the set computed in (ii). The latter is a single point,
which with (i) would give the containment excluded by (iii).

Now $R'$ is subdirect by (i) and (iv), it is closed under substitution in its
first $|A|$ coordinates, and $|B_{k-1}/\sigma_{k}|\le|A|$, so
Lemma~\ref{lem:full-fibre} supplies $d$ with
$(B_{k-1}/\sigma_{k})^{|A|}\times\{d\}\subseteq R'$. The associated set
$C^{\star}\subseteq B_{n-1}/\sigma_{n}$ is therefore nonempty, and proper by
(iii); so Lemma~\ref{lem:central-relation} applied to $R'^{-1}$ yields either a
proper nonempty central subuniverse of $B_{n-1}/\sigma_{n}$ or, through
Lemma~\ref{lem:power-to-base}, a proper nonempty binary absorbing subuniverse
of $B_{k-1}/\sigma_{k}$ --- both excluded. So subcase 1B ends at its first
branch.

\medskip\noindent\textbf{Case 2: $k>m$.} Then $B_{k}\subseteq B_{k-1}\subseteq
B_{m}\subseteq B_{m-1}$. The inductive hypothesis at $k-1$, with the same
$\delta$ and $m$, makes $((B_{k-1}\circ\delta)\cap B_{m-1})/\sigma_{m}$ either
$B_{m}/\sigma_{m}$ or $B_{m-1}/\sigma_{m}$.

If it is $B_{m}/\sigma_{m}$, that set is a single point, and
$\varnothing\neq B_{k}\subseteq(B_{k}\circ\delta)\cap B_{m-1}\subseteq
(B_{k-1}\circ\delta)\cap B_{m-1}$ forces
$((B_{k}\circ\delta)\cap B_{m-1})/\sigma_{m}=B_{m}/\sigma_{m}$ as well; the
additional clause is carried over unchanged from the inductive hypothesis,
which had the same $\delta$ and the same $m$. So assume
\begin{equation}\label{eq:sip-case2}
  ((B_{k-1}\circ\delta)\cap B_{m-1})/\sigma_{m}=B_{m-1}/\sigma_{m},
\end{equation}
and split on $T_{k}$. In each subcase the answer is $B_{m-1}/\sigma_{m}$, so
the additional clause is not triggered.

\emph{Subcase 2A: $T_{k}\in\{\TBA,\TC\}$.} The formula
$\exists y\,(X(y)\wedge\delta(x,y))\wedge B_{m-1}(x)$ defines
$(B_{k-1}\circ\delta)\cap B_{m-1}$ at $X=B_{k-1}$ and
$(B_{k}\circ\delta)\cap B_{m-1}$ at $X=B_{k}$, and the latter is nonempty as it
contains $B_{k}$; so Lemmas~\ref{lem:pp-propagation} and
\ref{lem:factor-type} give
$((B_{k}\circ\delta)\cap B_{m-1})/\sigma_{m}\subte{T_{k}}
((B_{k-1}\circ\delta)\cap B_{m-1})/\sigma_{m}=B_{m-1}/\sigma_{m}$. A
\emph{proper} such containment is a proper nonempty binary absorbing or central
subuniverse of $B_{m-1}/\sigma_{m}$, which is excluded, so it is an equality.

\emph{Subcase 2B: $T_{k}=\TS$.} Choose a nonempty $G\subseteq B_{k}$ with
$G\subt{\TBA,\TC}B_{k-1}$ and run 2A with $G$ for $B_{k}$: it gives
$((G\circ\delta)\cap B_{m-1})/\sigma_{m}=B_{m-1}/\sigma_{m}$, and
$G\subseteq B_{k}$ makes the left side a subset of
$((B_{k}\circ\delta)\cap B_{m-1})/\sigma_{m}$.

\emph{Subcase 2C: $T_{k}=\TD$.} Put
\[
  R=\bigl\{(a/\sigma_{k},b/\sigma_{m}):a\in B_{k-1},\ b\in B_{m-1},\
    (a,b)\in\delta\bigr\}\le\mathbf B_{k-1}/\sigma_{k}\times
    \mathbf B_{m-1}/\sigma_{m} .
\]
It is subdirect: $B_{k-1}\subseteq B_{m-1}$ lets $b=a$ serve for the left
projection, and the right projection is $B_{m-1}/\sigma_{m}$ by
\eqref{eq:sip-case2}. Moreover $B_{k-1}/\sigma_{k}\times B_{m}/\sigma_{m}
\subseteq R$: here $k>m$ gives $B_{k-1}\subseteq B_{m}$, so for $a\in B_{k-1}$
the pair $(a,a)$ witnesses $(a/\sigma_{k},a/\sigma_{m})\in R$ with
$a/\sigma_{m}$ the single class of $B_{m}/\sigma_{m}$. \textbf{This is the only
use of $k>m$ in the subcase, and without it $R$ need not have a centre at
all.}

So the set $C^{\star}=\{c\in B_{m-1}/\sigma_{m}:\{c\}\times B_{k-1}/\sigma_{k}
\subseteq R^{-1}\}$ is nonempty. Lemma~\ref{lem:central-relation} applied to
$R^{-1}\sdle\mathbf B_{m-1}/\sigma_{m}\times\mathbf B_{k-1}/\sigma_{k}$ leaves
only its alternative (b), alternative (c) being excluded by BA-freeness of
$B_{k-1}/\sigma_{k}$; and since $B_{m-1}/\sigma_{m}$ is centre free the central
subuniverse $C^{\star}$ cannot be proper. Hence
$C^{\star}=B_{m-1}/\sigma_{m}$ and $R$ is the full product.

Finally, fullness transfers from $B_{k-1}$ to $B_{k}$. Let $a\in B_{k}$ and
$b\in B_{m-1}$. Fullness gives $a'\in B_{k-1}$ and $b'\in B_{m-1}$ with
$(a',b')\in\delta$, $a'/\sigma_{k}=a/\sigma_{k}$ and
$b'/\sigma_{m}=b/\sigma_{m}$. From $a\in B_{k}\subseteq E_{k}$ and
$a'\,\sigma_{k}\,a$ we get $a'\in E_{k}$, so $a'\in B_{k-1}\cap E_{k}=B_{k}$;
hence $b'\in(B_{k}\circ\delta)\cap B_{m-1}$ and $b'/\sigma_{m}=b/\sigma_{m}$.
As $b$ was arbitrary, $((B_{k}\circ\delta)\cap B_{m-1})/\sigma_{m}=
B_{m-1}/\sigma_{m}$.
\end{proof}

\begin{hazard}[Three steps this proof supplies]\label{haz:self-intersection}
Subcase 1B's first branch is where the projection is actually computed, and it
needs two things beyond the displayed identity
$(E\cap B_{k-1})/\sigma_{k}=B_{k-1}/\sigma_{k}$: that a \emph{single}
$\delta$-class $E$ works, which comes from the witnesses for one tuple
enumerating all of $B_{k-1}/\sigma_{k}$ and so from $|B_{k-1}/\sigma_{k}|\le
|A|$ again; and that $E$ \emph{meets} $B_{k}$, without which $E$ has nothing to
do with $B_{k}\circ\delta$. The second follows from the first, since the
classes exhausted include the class of $B_{k}$.

Subcase 2C ends with a transfer step: $R$ full is a statement about
$B_{k-1}$, and the conclusion is about $B_{k}$. It goes through only because
$B_{k}$ is cut out of $B_{k-1}$ by a $\sigma_{k}$-block, so a $\sigma_{k}$-move
inside $B_{k-1}$ cannot leave $B_{k}$ --- the same observation that drives
subcase 1A.

The proof does not need $\delta$ to be irreducible, or comparable with any
$\sigma_{i}$; it is an arbitrary congruence, and the additional clause is the
only place any relation between $\delta$ and $\sigma_{m}$ is asserted.
\end{hazard}

\begin{longlemma}[The intersection property]\label{lem:intersection-good}
Let $B\lll A$, $C\lll A$ with $B\cap C\neq\varnothing$. Then
\begin{enumerate}\alphlabels\tightlist
\item[\textup{(d)}] if $\delta$ is a dividing congruence for $B\le A$ then
  $|(B\cap C)/\delta|=1$ or $(B\cap C)/\delta=B/\delta$; and if
  $|(B\cap C)/\delta|=1$ then $\delta\supseteq\delta_{1}\cap\dots\cap
  \delta_{s}$, where $\delta_{1},\dots,\delta_{s}$ are the congruences of the
  \emph{chosen} chains witnessing $B\lll A$ and $C\lll A$;
\item[\textup{(s)}] if $G\subt{\TBA,\TC}B$ then $G\cap C\neq\varnothing$.
\end{enumerate}
\end{longlemma}

\begin{proof}
Fix chains
\[
  B=B_{k}\subt[A]{T_{k}(\sigma_{k})}\dots\subt[A]{T_{1}(\sigma_{1})}B_{0}=A,
  \qquad
  C=C_{\ell}\subt[A]{\mathcal T_{\ell}(\omega_{\ell})}\dots
  \subt[A]{\mathcal T_{1}(\omega_{1})}C_{0}=A
\]
witnessing $B\lll A$ and $C\lll A$, with all $T_{i},\mathcal T_{j}\in
\{\TBA,\TC,\TS,\TD\}$ and, by Definition~\ref{def:types}, $\sigma_{i}=A^{2}$
whenever $T_{i}\neq\TD$ and $\omega_{j}=A^{2}$ whenever
$\mathcal T_{j}\neq\TD$. Put $\sigma=\bigcap_{i=1}^{k}\sigma_{i}$ and
$\omega=\bigcap_{j=1}^{\ell}\omega_{j}$, so that
$\sigma\cap\omega=\delta_{1}\cap\dots\cap\delta_{s}$ is exactly the
intersection named in the ``moreover'' clause. The induction is on $k+\ell$
over \emph{pairs of chains} --- every appeal below names the shorter pair it is
made to --- and proves (s) and (d) together.

Three facts are used throughout and are recorded once.

\emph{(F1) $B^{2}\subseteq\sigma$ and $C^{2}\subseteq\omega$.} If $T_{i}=\TD$
then $B_{i}=B_{i-1}\cap E_{i}$ for a block $E_{i}$ of $\sigma_{i}$, so
$B_{i}^{2}\subseteq\sigma_{i}$ and a fortiori $B^{2}\subseteq\sigma_{i}$;
otherwise $\sigma_{i}=A^{2}$. Likewise for $C$. In particular
$(a,b)\in\sigma\cap\omega$ for all $a,b\in B\cap C$.

\emph{(F2) $\sigma\subseteq\sigma_{i}$ and $\omega\subseteq\omega_{j}$} for
every $i$ and $j$, by construction.

\emph{(F3) $|A|\ge2$.} That $\delta$ is a dividing congruence for $B\le A$
means, by Definition~\ref{def:types}, that $B'\subt[A]{\TD(\delta)}B$ for some
$B'$; such a $B'$ is a nonempty proper subset $B\cap E$ of $B$ cut out by a
$\delta$-block, so $|B/\delta|\ge2$, and hence $|A|\ge|B|\ge2$. The same
appeal supplies the two properties of $\delta$ used below: $\mathbf B/\delta$
is BA and centre free, and $B^{2}\subseteq\cover{\delta}$.

\medskip\noindent\textbf{Part (s).} Let $G\subt{\TBA,\TC}B$. If $\ell=0$ then
$C=A\supseteq G\neq\varnothing$. So let $\ell\ge1$. The chain for $C$
truncated at $C_{\ell-1}$ witnesses $C_{\ell-1}\lll A$ with $\ell-1$ steps, and
$B\cap C_{\ell-1}\supseteq B\cap C\neq\varnothing$; so the inductive hypothesis
(s), at $k+(\ell-1)$, gives
\begin{equation}\label{eq:ig-s-start}
  G\cap C_{\ell-1}\neq\varnothing .
\end{equation}
Write $A'=B\cap C_{\ell-1}$, a nonempty subuniverse of $\mathbf A$. Since
$G\subseteq B$ we have $G\cap X=G\cap A'\cap X$ for every $X\subseteq
C_{\ell-1}$, and $B\cap C_{\ell}\supseteq B\cap C\neq\varnothing$. By
Lemma~\ref{lem:ba-center-intersect} applied twice, once for $\TBA$ and once for
$\TC$,
\begin{equation}\label{eq:ig-s-cut}
  G\cap C_{\ell-1}=G\cap A'\dsubte{\TBA,\TC}A',
\end{equation}
and by \eqref{eq:ig-s-start} the dot may be dropped. Split on
$\mathcal T_{\ell}$.

\emph{Case $\mathcal T_{\ell}=\TD$,} with $C_{\ell}=C_{\ell-1}\cap E$ for a
block $E$ of $\omega_{\ell}$. Apply the inductive hypothesis (d) to the pair
$(C_{\ell-1},B)$ --- that is, with $C_{\ell-1}$ in the role of the first
subuniverse and $B$ in the role of the second --- and to the dividing
congruence $\omega_{\ell}$ for $C_{\ell-1}\le A$, at $(\ell-1)+k$. Either
$|A'/\omega_{\ell}|=1$ or $A'/\omega_{\ell}=C_{\ell-1}/\omega_{\ell}$.

In the first case the sole $\omega_{\ell}$-class meeting $A'$ is $E$, because
$\varnothing\neq B\cap C\subseteq A'\cap E$; hence $A'\subseteq E$ and
$A'=A'\cap E=B\cap C_{\ell}$, so
$G\cap C_{\ell}=G\cap A'\cap C_{\ell}=G\cap A'=G\cap C_{\ell-1}\neq\varnothing$.

In the second case Lemma~\ref{lem:factor-type} applied to
\eqref{eq:ig-s-cut} gives
$(G\cap C_{\ell-1})/\omega_{\ell}\subte{\TBA,\TC}A'/\omega_{\ell}=
C_{\ell-1}/\omega_{\ell}$. If that containment is proper it exhibits a proper
nonempty binary absorbing subuniverse of $C_{\ell-1}/\omega_{\ell}$,
contradicting the requirement of Definition~\ref{def:types}(iii) for the step
$C_{\ell}\subt[A]{\TD(\omega_{\ell})}C_{\ell-1}$ that
$\mathbf C_{\ell-1}/\omega_{\ell}$ be BA and centre free. If it is an equality
then every $\omega_{\ell}$-class meeting $C_{\ell-1}$ already meets
$G\cap C_{\ell-1}$; the class $E$ does meet $C_{\ell-1}$, since
$C_{\ell}\neq\varnothing$, so there is
$g\in G\cap C_{\ell-1}\cap E=G\cap C_{\ell}$.

\emph{Case $\mathcal T_{\ell}\in\{\TBA,\TC\}$.} By
Lemma~\ref{lem:ba-center-intersect}, $B\cap C_{\ell}=A'\cap C_{\ell}
\dsubte{\mathcal T_{\ell}}A'$, and it is nonempty. If
$G\cap C_{\ell-1}=A'$ then $A'\subseteq G$, and $\varnothing\neq B\cap
C_{\ell}\subseteq A'$ gives $G\cap C_{\ell}\neq\varnothing$. If
$B\cap C_{\ell}=A'$ then $G\cap C_{\ell}=G\cap A'\cap C_{\ell}=G\cap A'\neq
\varnothing$. Otherwise both containments are proper, so
$G\cap C_{\ell-1}\subt{\TBA,\TC}A'$ and $B\cap C_{\ell}
\subt{\mathcal T_{\ell}}A'$, and Lemma~\ref{lem:ba-central-meets} inside
$\mathbf A'$ makes their intersection nonempty; that intersection is contained
in $G\cap C_{\ell}$.

\emph{Case $\mathcal T_{\ell}=\TS$.} By Definition~\ref{def:ba-center-free}
there is a nonempty $G'\subseteq C_{\ell}$ with $G'\subt{\TBA,\TC}C_{\ell-1}$
--- proper, since $G'\subseteq C_{\ell}\lneq C_{\ell-1}$. The inductive
hypothesis (s) applied to the pair $(C_{\ell-1},B)$ with the subuniverse $G'$,
at $(\ell-1)+k$, gives $B\cap G'\neq\varnothing$; so
$B\cap G'=A'\cap G'\subte{\TBA,\TC}A'$ by
Lemma~\ref{lem:ba-center-intersect}. If $G\cap C_{\ell-1}=A'$ then
$\varnothing\neq B\cap G'\subseteq A'\subseteq G$ and $B\cap G'\subseteq
C_{\ell}$, so $G\cap C_{\ell}\neq\varnothing$. If $B\cap G'=A'$ then
$A'\subseteq G'\subseteq C_{\ell}$, so $G\cap C_{\ell}\supseteq G\cap A'\neq
\varnothing$. Otherwise both are proper and
Lemma~\ref{lem:ba-central-meets} gives
$\varnothing\neq(G\cap C_{\ell-1})\cap(B\cap G')\subseteq G\cap C_{\ell}$.

\medskip\noindent\textbf{Part (d).} For $0\le m\le k$ and $0\le n\le\ell$ put
\[
  P_{m,n}=\bigl\{(a_{1},\dots,a_{|A|})\in(B_{m}\cap C_{n})^{|A|}:
    (a_{i},a_{j})\in\sigma\cap\omega\ \text{for all }i,j\bigr\},
  \qquad S_{m,n}=P_{m,n}/\delta ,
\]
the quotient being coordinatewise and, as always, an \emph{image}
(\S\ref{sec:legislation}, item~\ref{leg:quotient}): a tuple
$x\in(A/\delta)^{|A|}$ lies in $S_{m,n}$ exactly when \emph{one} tuple of
representatives of its entries lies in $P_{m,n}$. Each $P_{m,n}$ is a
subuniverse of $\mathbf A^{|A|}$, being cut out by unary and binary
conditions, so each $S_{m,n}$ is a subuniverse of $(\mathbf A/\delta)^{|A|}$;
each is closed under substitution $x\mapsto x\circ g$ along any
$g:[|A|]\to[|A|]$, because every defining condition is unary in $a_{i}$ or
binary in $(a_{i},a_{j})$, and in particular each is totally symmetric and
shift-closed; and $S_{m,n}\supseteq S_{m',n'}$ whenever $m\le m'$ and
$n\le n'$. Write $S=S_{0,0}$; since $B_{0}\cap C_{0}=A$ and
$(a,a)\in\sigma\cap\omega$ for every $a$, this one --- and only this one --- is
also reflexive.

By (F3) the arity $|A|$ is at least $2$, so Lemma~\ref{lem:tot-sym-irred}
applies to $S$ and the dividing congruence $\delta$ for $B\le A$, leaving two
cases.

\emph{Case 1: $S=\{(a/\delta,\dots,a/\delta):a\in A\}$.} Let
$(a,b)\in\sigma\cap\omega$. Since $|A|\ge2$ the tuple $(a,b,b,\dots,b)$ lies
in $P_{0,0}$, its entries being pairwise $\sigma\cap\omega$-related; so its
image lies in $S$ and is therefore constant, giving $a/\delta=b/\delta$. Hence
$\sigma\cap\omega\subseteq\delta$, which is the ``moreover'' clause. By (F1)
any $a,b\in B\cap C$ satisfy $(a,b)\in\sigma\cap\omega\subseteq\delta$, and
$B\cap C\neq\varnothing$, so $|(B\cap C)/\delta|=1$.

\emph{Case 2: $(B/\delta)^{|A|}\subseteq S$.} We show that the only surviving
subcase is 2C, in which $(B\cap C)/\delta=B/\delta$; by (F3) that alternative
has $|(B\cap C)/\delta|=|B/\delta|\ge2$, so the ``moreover'' clause is not
triggered in Case 2 and nothing further is owed.

\emph{Subcase 2A: $(B/\delta)^{|A|}\not\subseteq S_{k,0}$.} Choose $m$ minimal
with $(B/\delta)^{|A|}\not\subseteq S_{m,0}$; since $S_{0,0}=S$ we have
$m\ge1$ and $(B/\delta)^{|A|}\subseteq S_{m-1,0}$, so $(B/\delta)^{|A|}$ is a
subalgebra of $\mathbf S_{m-1,0}$. Note that for every $b\in B$ the constant
tuple $(b,\dots,b)$ lies in $P_{m,0}$, because $B\subseteq B_{m}$ and
$(b,b)\in\sigma\cap\omega$; so
\begin{equation}\label{eq:ig-2A-nonempty}
  \varnothing\neq S_{m,0}\cap(B/\delta)^{|A|}\subsetneq(B/\delta)^{|A|},
\end{equation}
the properness being the choice of $m$. Split on $T_{m}$.

\emph{Subsubcase 2A1: $T_{m}\in\{\TBA,\TC\}$.} The formula
\[
  \Phi_{X}(x_{1},\dots,x_{|A|})=
  \bigwedge_{i,j}(\sigma\cap\omega)(x_{i},x_{j})\wedge\bigwedge_{i}X(x_{i})
\]
is positive primitive in the unary relation $X$, and defines $P_{m-1,0}$ at
$X=B_{m-1}$ and $P_{m,0}$ at $X=B_{m}$. Since $B_{m}\subt{T_{m}}B_{m-1}$,
Lemma~\ref{lem:pp-propagation} gives $P_{m,0}\dsubte{T_{m}}P_{m-1,0}$, and
$P_{m,0}\neq\varnothing$ by the constant tuples, so
$P_{m,0}\subte{T_{m}}P_{m-1,0}$. Lemma~\ref{lem:factor-type}, applied to that
containment inside $\mathbf P_{m-1,0}$ with the congruence induced by
$\delta^{|A|}$, gives $S_{m,0}\subte{T_{m}}S_{m-1,0}$, and then
Lemma~\ref{lem:ba-center-intersect} restricts along the subalgebra
$(B/\delta)^{|A|}\le\mathbf S_{m-1,0}$:
\[
  S_{m,0}\cap(B/\delta)^{|A|}\dsubte{T_{m}}(B/\delta)^{|A|} .
\]
By \eqref{eq:ig-2A-nonempty} the left side is nonempty and proper, so this is
$\subt{T_{m}}$, and Lemma~\ref{lem:power-to-base} produces a proper nonempty
subuniverse of $\mathbf B/\delta$ of type $T_{m}\in\{\TBA,\TC\}$. That
contradicts (F3).

\emph{Subsubcase 2A2: $T_{m}=\TS$.} By Definition~\ref{def:ba-center-free}
choose a nonempty $G\subseteq B_{m}$ with $G\subt{\TBA,\TC}B_{m-1}$. Since
$B\lll^{A}B_{m-1}\lll A$ along the two halves of the chain,
Lemma~\ref{lem:strong-nonempty} gives $G\cap B\neq\varnothing$. Run 2A1 with
$G$ in place of $B_{m}$: writing $P^{G}=\Phi_{G}$ and $S^{G}=P^{G}/\delta$,
the same three lemmas give $S^{G}\cap(B/\delta)^{|A|}\dsubte{\TBA}
(B/\delta)^{|A|}$; it is nonempty, since $b\in G\cap B$ contributes its
constant tuple, and proper, since $G\subseteq B_{m}$ makes it a subset of the
proper set in \eqref{eq:ig-2A-nonempty}. Lemma~\ref{lem:power-to-base} again
contradicts (F3). (Only the $\TBA$ half of $G\subt{\TBA,\TC}B_{m-1}$ is spent
here.)

\emph{Subsubcase 2A3: $T_{m}=\TD$,} with $B_{m}=B_{m-1}\cap E$ for a block $E$
of $\sigma_{m}$. Let $B\circ\delta$ denote the $\delta$-saturation of $B$,
a subuniverse of $\mathbf A$, and note $a/\delta\in B/\delta\iff a\in
B\circ\delta$. Put
\[
  P'=\bigl\{(a_{1},\dots,a_{|A|},b): b\in B_{m-1},\ \forall i\;
    a_{i}\in B_{m-1}\cap(B\circ\delta),\ \forall i,j\;(a_{i},a_{j})\in
    \sigma\cap\omega,\ \forall i\;(a_{i},b)\in\sigma_{m}\bigr\}
\]
and let $R'\le(\mathbf A/\delta)^{|A|}\times\mathbf A/\sigma_{m}$ be its image
under $(a_{1},\dots,a_{|A|},b)\mapsto(a_{1}/\delta,\dots,a_{|A|}/\delta,
b/\sigma_{m})$, so that $R'\subseteq(B/\delta)^{|A|}\times B_{m-1}/\sigma_{m}$.
Four facts about $R'$.

\emph{(i) $\proj_{1,\dots,|A|}(R')=(B/\delta)^{|A|}$.} Let
$x\in(B/\delta)^{|A|}$. By minimality of $m$, $x\in S_{m-1,0}$, so there are
$a_{i}\in B_{m-1}$, pairwise $\sigma\cap\omega$-related, with
$a_{i}/\delta=x_{i}$; each $a_{i}$ then lies in $B\circ\delta$. By (F2),
$(a_{i},a_{j})\in\sigma\subseteq\sigma_{m}$, so $b=a_{1}\in B_{m-1}$ satisfies
$(a_{i},b)\in\sigma_{m}$ for every $i$, and $(a_{1},\dots,a_{|A|},b)\in P'$.

\emph{(ii) $\proj_{|A|+1}(R')=((B\circ\delta)\cap B_{m-1})/\sigma_{m}$.} For
$\subseteq$: a witnessing tuple has $a_{1}\in B_{m-1}\cap(B\circ\delta)$ and
$(a_{1},b)\in\sigma_{m}$, so $b/\sigma_{m}=a_{1}/\sigma_{m}$. For
$\supseteq$: given $c$ in that set, the tuple $(c,\dots,c,c)$ lies in $P'$.

\emph{(iii) $(B/\delta)^{|A|}\times B_{m}/\sigma_{m}\not\subseteq R'$.} Suppose
otherwise and fix $b_{0}\in B\subseteq B_{m}$. For any
$x\in(B/\delta)^{|A|}$ the pair $(x,b_{0}/\sigma_{m})$ lies in $R'$, so it has
a witness $(a_{1},\dots,a_{|A|},b)\in P'$ with $b/\sigma_{m}=b_{0}/\sigma_{m}$;
then $(a_{i},b_{0})\in\sigma_{m}$, and $b_{0}\in B_{m}\subseteq E$ forces
$a_{i}\in E$, whence $a_{i}\in B_{m-1}\cap E=B_{m}$. So the witness lies in
$P_{m,0}$ and $x\in S_{m,0}$ --- for every $x$, contradicting the choice of
$m$.

\emph{(iv) $\proj_{|A|+1}(R')=B_{m-1}/\sigma_{m}$.} Apply
Lemma~\ref{lem:self-intersection-pc} to the chain for $B$, the congruence
$\delta$ and the index $m$: by (ii) the projection is either
$B_{m-1}/\sigma_{m}$ or $B_{m}/\sigma_{m}$. In the second case
$B_{m}\subseteq E$ makes $B_{m}/\sigma_{m}$ a single point, so by (i) every
$x\in(B/\delta)^{|A|}$ is paired in $R'$ with that point, i.e.\
$(B/\delta)^{|A|}\times B_{m}/\sigma_{m}\subseteq R'$, contradicting (iii).

By (i) and (iv), $R'$ is subdirect in $(\mathbf B/\delta)^{|A|}\times
\mathbf B_{m-1}/\sigma_{m}$; both factors are finite algebras with the single
basic operation $w$, so Convention~\ref{conv:standing} holds for them. $R'$ is
closed under substitution in its first $|A|$ coordinates, since every defining
condition of $P'$ is unary in $a_{i}$ or binary in $(a_{i},a_{j})$ or
$(a_{i},b)$; and $|B/\delta|\le|A|$. So Lemma~\ref{lem:full-fibre}, with
$F=B/\delta$ and $N=|A|$, supplies $d\in B_{m-1}/\sigma_{m}$ with
$(B/\delta)^{|A|}\times\{d\}\subseteq R'$.

Now apply Lemma~\ref{lem:central-relation} to the transpose
$R'^{-1}=\{(y,x):(x,y)\in R'\}\sdle\mathbf B_{m-1}/\sigma_{m}\times
(\mathbf B/\delta)^{|A|}$, whose associated set is
$C^{\star}=\{c\in B_{m-1}/\sigma_{m}:\{c\}\times(B/\delta)^{|A|}\subseteq
R'^{-1}\}$. Alternative (a) fails, since $d\in C^{\star}$. And $C^{\star}$ is
\emph{proper}: were it all of $B_{m-1}/\sigma_{m}$ we should have
$R'=(B/\delta)^{|A|}\times B_{m-1}/\sigma_{m}$, which contains
$(B/\delta)^{|A|}\times B_{m}/\sigma_{m}$ and contradicts (iii). So
alternative (b) makes $C^{\star}$ a proper nonempty central subuniverse of
$B_{m-1}/\sigma_{m}$, contradicting Definition~\ref{def:types}(iii) for the
step $B_{m}\subt[A]{\TD(\sigma_{m})}B_{m-1}$; and alternative (c) gives a
proper nonempty binary absorbing subuniverse of the \emph{power}
$(B/\delta)^{|A|}$, which Lemma~\ref{lem:power-to-base} moves down to
$\mathbf B/\delta$, contradicting (F3). Both obligations of
Caveat~\ref{haz:central-relation} are thereby discharged.

\emph{Subcase 2B: $(B/\delta)^{|A|}\subseteq S_{k,0}$ but
$(B/\delta)^{|A|}\not\subseteq S_{k,\ell}$.} Choose $n$ minimal with
$(B/\delta)^{|A|}\not\subseteq S_{k,n}$; then $n\ge1$ and
$(B/\delta)^{|A|}\subseteq S_{k,n-1}$. Fixing once and for all $b_{1}\in
B\cap C\subseteq B_{k}\cap C_{n}$ --- this is the one place where the
hypothesis $B\cap C\neq\varnothing$ is spent inside Case 2 --- its constant
tuple lies in $P_{k,n}$, so
\begin{equation}\label{eq:ig-2B-nonempty}
  \varnothing\neq S_{k,n}\cap(B/\delta)^{|A|}\subsetneq(B/\delta)^{|A|} .
\end{equation}
Split on $\mathcal T_{n}$.

\emph{Subsubcase 2B1: $\mathcal T_{n}\in\{\TBA,\TC\}$.} Verbatim 2A1, with the
pp-formula
$\bigwedge_{i,j}(\sigma\cap\omega)(x_{i},x_{j})\wedge\bigwedge_{i}B_{k}(x_{i})
\wedge\bigwedge_{i}X(x_{i})$, with $X$ instantiated at $C_{n-1}$ and at
$C_{n}$, and with \eqref{eq:ig-2B-nonempty} in place of
\eqref{eq:ig-2A-nonempty}. It contradicts (F3).

\emph{Subsubcase 2B2: $\mathcal T_{n}=\TS$.} Choose a nonempty $G\subseteq
C_{n}$ with $G\subt{\TBA,\TC}C_{n-1}$. Here $B$ and $C_{n-1}$ lie on different
chains, so Lemma~\ref{lem:strong-nonempty} does not apply and the induction is
used instead: the inductive hypothesis (s) for the pair $(C_{n-1},B)$, at
$(n-1)+k$, with $C_{n-1}\cap B\supseteq C\cap B\neq\varnothing$, gives
$G\cap B\neq\varnothing$. Now run 2A2 with $G$ in place of $C_{n}$: the
witness $b\in G\cap B$ lies in $B_{k}\cap C_{n-1}$ and makes the intersection
nonempty, $G\subseteq C_{n}$ makes it proper by
\eqref{eq:ig-2B-nonempty}, and Lemma~\ref{lem:power-to-base} contradicts (F3).

\emph{Subsubcase 2B3: $\mathcal T_{n}=\TD$,} with $C_{n}=C_{n-1}\cap E$ for a
block $E$ of $\omega_{n}$. Put
\[
  P'=\bigl\{(a_{1},\dots,a_{|A|},b): b\in C_{n-1},\ \forall i\;
    a_{i}\in B_{k}\cap C_{n-1},\ \forall i,j\;(a_{i},a_{j})\in\sigma\cap\omega,
    \ \forall i\;(a_{i},b)\in\omega_{n}\bigr\}
\]
and let $R'$ be its image in $(\mathbf A/\delta)^{|A|}\times\mathbf A/
\omega_{n}$. No saturation appears in the first $|A|$ coordinates this time,
because $a_{i}\in B_{k}=B$ already gives $a_{i}/\delta\in B/\delta$.

\emph{(i)} $\proj_{1,\dots,|A|}(R')=(B/\delta)^{|A|}$: for
$x\in(B/\delta)^{|A|}\subseteq S_{k,n-1}$ take witnesses $a_{i}\in B_{k}\cap
C_{n-1}$; by (F2) $(a_{i},a_{j})\in\omega\subseteq\omega_{n}$, so
$b=a_{1}\in C_{n-1}$ serves.
\emph{(ii)} $\proj_{|A|+1}(R')=(B_{k}\cap C_{n-1})/\omega_{n}$, by the
argument of 2A3(ii).
\emph{(iii)} $(B/\delta)^{|A|}\times C_{n}/\omega_{n}\not\subseteq R'$: with
$b_{1}\in B\cap C\subseteq B_{k}\cap C_{n}$ in place of $b_{0}$, the argument
of 2A3(iii) puts every witness in $C_{n-1}\cap E=C_{n}$ and so every
$x\in(B/\delta)^{|A|}$ in $S_{k,n}$, contradicting the choice of $n$.
\emph{(iv)} $\proj_{|A|+1}(R')=C_{n-1}/\omega_{n}$: the inductive hypothesis
(d) for the pair $(C_{n-1},B_{k})$ with the dividing congruence $\omega_{n}$
for $C_{n-1}\le A$, at $(n-1)+k$, leaves $|(B_{k}\cap C_{n-1})/\omega_{n}|=1$
or $(B_{k}\cap C_{n-1})/\omega_{n}=C_{n-1}/\omega_{n}$. In the first case the
sole class is $E$, since $b_{1}\in B_{k}\cap C_{n}\subseteq E$; so
$B_{k}\cap C_{n-1}\subseteq E$ and $B_{k}\cap C_{n-1}=B_{k}\cap C_{n-1}\cap
E=B_{k}\cap C_{n}$, whence $P_{k,n-1}=P_{k,n}$ and $S_{k,n-1}=S_{k,n}$,
contradicting the choice of $n$.

From here 2A3 runs word for word with $\omega_{n}$ for $\sigma_{m}$,
$C_{n-1}$ for $B_{m-1}$ and $C_{n}$ for $B_{m}$: a full fibre by
Lemma~\ref{lem:full-fibre}, a proper
nonempty $C^{\star}$, and Lemma~\ref{lem:central-relation} contradicting either
Definition~\ref{def:types}(iii) for $C_{n}\subt[A]{\TD(\omega_{n})}C_{n-1}$ or
(F3).

\emph{Subcase 2C: $(B/\delta)^{|A|}\subseteq S_{k,\ell}$.} For $b\in B$ the
constant tuple $(b/\delta,\dots,b/\delta)$ lies in $S_{k,\ell}$, so it has a
witness in $(B_{k}\cap C_{\ell})^{|A|}=(B\cap C)^{|A|}$, whose first entry $a$
satisfies $a/\delta=b/\delta$ with $a\in B\cap C$. Hence $B/\delta\subseteq
(B\cap C)/\delta$, and the reverse containment is trivial, so
$(B\cap C)/\delta=B/\delta$.
\end{proof}

\begin{hazard}[The four places this proof can be misread]\label{haz:intersection-draft}
\begin{enumerate}\alphlabels\tightlist
\item \emph{$S_{m,n}$ is an image, not an intersection.} A tuple lies in it
  when \emph{one} tuple of representatives satisfies all the conditions at
  once, and those conditions couple the representatives to each other. So
  $S_{m,n}$ is not a power of $(B_{m}\cap C_{n})/\delta$ cut down by $S$, and
  the membership tests in 2A3(iii), 2B3(iii) and 2C all work by producing or
  constraining a single witnessing tuple.
\item \emph{Lemmas~\ref{lem:ba-center-intersect} and \ref{lem:factor-type}
  deliver $\subte{}$, never $\subt{}$.} Part (s) therefore splits five times on
  whether a containment is proper, and in each case the equality branch
  \emph{proves the lemma} rather than contradicting anything: an absorbing
  subuniverse that is everything is no obstruction, but it does make the set on
  the left swallow the set one is trying to meet. Reading these lemmas as
  producing proper subuniverses collapses all five splits, and it is the most
  tempting error available here.
\item \emph{The second alternative of Lemma~\ref{lem:self-intersection-pc} has
  to be excluded by hand}, in 2A3(iv), and likewise its analogue in 2B3(iv).
  What kills it is that $B_{m}$ lies inside a single $\sigma_{m}$-block: the
  alternative would make the last coordinate of $R'$ constant and hence, with
  surjectivity of the first $|A|$, force exactly the containment that the
  minimal choice of $m$ forbids. In 2B3(iv) the same manoeuvre also identifies
  \emph{which} class the single one is --- the one defining $C_{n}$, witnessed
  by a point of $B\cap C$.
\item \emph{Every appeal to Lemma~\ref{lem:central-relation} owes two
  discharges}, per Caveat~\ref{haz:central-relation}: the centre must be shown
  proper --- here by the same non-containment that the choice of $m$ or $n$
  provides --- and the binary absorbing subuniverse of alternative (c) lives on
  a power and must be moved down by Lemma~\ref{lem:power-to-base}.
\end{enumerate}

Two dependencies remain outside the proof: it rests on
Lemma~\ref{lem:self-intersection-pc}, which is still an outline, and on the
imported Lemma~\ref{lem:power-to-base}. Minimality of the chains for $B$ and
$C$ is \emph{not} among them, and could not be: the inductive appeals in part
(s), 2B2 and 2B3 are to truncations of the chain for $C$, which need not be
minimal even when the original is.
\end{hazard}

\begin{corollary}\label{cor:intersection-good}
Let $R\sdle\mathbf A_{1}\times\mathbf A_{2}$, $B_{1}\lll A_{1}$,
$B_{2}\lll A_{2}$, and let $\sigma$ be a dividing congruence for
$B_{1}\lll A_{1}$. Then $\proj_{1}\bigl(R\cap(B_{1}\times B_{2})\bigr)/\sigma$
is empty, of size $1$, or equal to $B_{1}/\sigma$.
\end{corollary}

\begin{proof}
Let $f_{i}:\mathbf R\to\mathbf A_{i}$ be the projections. By
Lemma~\ref{lem:reverse-hom}(b), $f_{i}^{-1}(B_{i})\lll\mathbf R$; apply
Lemma~\ref{lem:intersection-good}(d) inside $\mathbf R$ to these two, with the
congruence $f_{1}^{-1}(\sigma)$, and read the conclusion back through $f_{1}$.
\end{proof}

\begin{hazard}\label{haz:cor-intersection}
The alternative ``empty'' is present in
Corollary~\ref{cor:intersection-good} and absent from
Lemma~\ref{lem:intersection-good}, which hypothesises $B\cap C\neq\varnothing$.
It is a real alternative, and every call site below has to exclude it
separately.
\end{hazard}

\begin{lemma}[Parallelogram property for $\TD$]\label{lem:parallelogram-D}
Let $B_{1},B_{2}\lll A$, $C_{1},C_{1}'\subt[A]{\TD(\sigma_{1})}B_{1}$ and
$C_{2},C_{2}'\subt[A]{\TD(\sigma_{2})}B_{2}$ with $C_{1}'\cap C_{2}$,
$C_{1}\cap C_{2}'$ and $C_{1}'\cap C_{2}'$ all nonempty. Then
$C_{1}\cap C_{2}\neq\varnothing$.
\end{lemma}

\begin{proof}
By Lemma~\ref{lem:intersection-good}(d), $(B_{1}\cap B_{2})/\sigma_{1}$ is of
size $1$ or all of $B_{1}/\sigma_{1}$. In the first case the single block is
the one defining $C_{1}$, because $C_{1}\cap C_{2}'\neq\varnothing$ puts a
point of $C_{1}$ in $B_{1}\cap B_{2}$; hence $B_{1}\cap B_{2}\subseteq C_{1}$
and $C_{1}\cap C_{2}=B_{1}\cap C_{2}\supseteq C_{1}'\cap C_{2}\neq\varnothing$.
The symmetric argument applies if $(B_{1}\cap B_{2})/\sigma_{2}$ has size $1$.

So assume $(B_{1}\cap B_{2})/\sigma_{1}=B_{1}/\sigma_{1}$ and
$(B_{1}\cap B_{2})/\sigma_{2}=B_{2}/\sigma_{2}$, and put
$S=\{(a/\sigma_{1},a/\sigma_{2}):a\in B_{1}\cap B_{2}\}$, which is then
subdirect in $(B_{1}/\sigma_{1})\times(B_{2}/\sigma_{2})$. Apply
Lemma~\ref{lem:intersection-good}(d) to each $\sigma_{1}$-class $C$ of $B_{1}$
against $B_{2}$ with the congruence $\sigma_{2}$; the hypothesis
$C\cap B_{2}\neq\varnothing$ holds because $S$ is subdirect. Each fibre of $S$
is therefore a single point or the whole of $B_{2}/\sigma_{2}$.

If some fibre is full, then either $S$ is full --- and then the fibre of
$C_{1}/\sigma_{1}$ contains $C_{2}/\sigma_{2}$, giving $C_{1}\cap C_{2}\neq
\varnothing$ --- or the centre of $S$ is proper and nonempty, and
Lemma~\ref{lem:central-relation} produces a proper nonempty binary absorbing or
central subuniverse of $B_{1}/\sigma_{1}$ or $B_{2}/\sigma_{2}$, contradicting
Definition~\ref{def:types}(iii). If every fibre is a single point, then the
fibre of $C_{1}'/\sigma_{1}$ is both $C_{2}/\sigma_{2}$ (witnessed by
$C_{1}'\cap C_{2}$) and $C_{2}'/\sigma_{2}$ (witnessed by $C_{1}'\cap C_{2}'$),
so $C_{2}=C_{2}'$ and $C_{1}\cap C_{2}=C_{1}\cap C_{2}'\neq\varnothing$.
\end{proof}

\subsection{Bridges}\label{sec:bridge-proofs}

Three of the statements below --- Lemmas~\ref{lem:nice-bridge-abelian},
\ref{lem:block-good-bridge} and \ref{lem:nontrivial-bridge} --- require their
bridge to be \emph{pair-reflexive}, and the hypothesis is not cosmetic in any
of them: Remark~\ref{rem:phi-twist} exhibits a bridge satisfying every other
hypothesis of the first and none of its conclusion. Lemma~\ref{lem:symmetrise}
is what makes the hypothesis free, since every bridge these statements are
applied to below is a $\delta^{\circ}$.

\begin{lemma}\label{lem:nice-bridge-abelian}
Let $\sigma$ be a congruence on $\mathbf A$ and $\delta$ a symmetric,
\emph{pair-reflexive} bridge from $\sigma$ to $\sigma$ with
$\proj_{1,2}(\delta)=\bcol\delta=A^{2}$. Then there is an abelian group
$(G;+,-)$ with $(A/\sigma;\delta/\sigma)\cong(G;x_{1}-x_{2}=x_{3}-x_{4})$.
\end{lemma}

\begin{proof}
Factor by $\sigma$ (Lemma~\ref{lem:irred-quotient}) and assume
$\sigma=0_{\mathbf A}$. By Lemma~\ref{lem:bridge-abelian} and
Lemma~\ref{lem:hobby-mckenzie}, $\mathbf A$ is affine, hence polynomially
equivalent to a $\mathbb Z_{p}$-module $\mathbf G$ and Maltsev via $x-y+z$.

View $\delta$ as a binary relation on $\mathbf A^{2}$ as in the proof of
Lemma~\ref{lem:bridge-abelian}. It is reflexive, by pair-reflexivity together
with $\proj_{1,2}(\delta)=A^{2}$, and symmetric by hypothesis. In a Maltsev
algebra a reflexive compatible binary relation is a congruence: if $\eta$ is
reflexive and compatible then for $(x,y),(y,z)\in\eta$ the Maltsev term applied
to $(x,y),(y,y),(y,z)$ gives $(x,z)\in\eta$. So $\delta$ is a congruence of
$\mathbf G^{2}$.

Clause (d) of Definition~\ref{def:bridge} says that a pair
$\bigl((x_{1},x_{2}),(x_{3},x_{4})\bigr)$ of $\delta$ has $x_{1}=x_{2}$ exactly
when $x_{3}=x_{4}$, so the diagonal of $\mathbf G^{2}$ is a $\delta$-block. A
congruence of $\mathbf G^{2}$ with the diagonal as a block is the kernel of
$(x,y)\mapsto x-y$ composed with a quotient of $\mathbf G$; since the diagonal
is exactly one block, no further collapsing occurs, and
$\delta=\{x_{1}-x_{2}=x_{3}-x_{4}\}$. Taking $G=\mathbf G$ finishes.
\end{proof}

\begin{remark}[Pair-reflexivity is not removable]\label{rem:phi-twist}
Without pair-reflexivity the correct conclusion is
$\delta=\{(x_{1},x_{2},x_{3},x_{4}):\varphi(x_{1}-x_{2})=x_{3}-x_{4}\}$ for an
automorphism $\varphi$ of $G$: factoring out the diagonal exhibits $\delta$ as
the preimage of a subgroup $K\le G\times G$ that meets $G\times 0$ and
$0\times G$ trivially and is subdirect, hence a graph. Symmetry says
$\varphi^{2}=\mathrm{id}$; pair-reflexivity says $\varphi=\mathrm{id}$.

Both $\varphi$ occur. On $\mathbf Z_{3}\in\mathcal V_{4}$ --- the operation
$x_{1}+x_{2}+x_{3}+x_{4}$ is idempotent since $4\equiv1$, and is a special WNU
--- take $\sigma=0$ and
\[
  \delta=\{(x_{1},x_{2},x_{3},x_{4})\in\mathbb Z_{3}^{4}:
    x_{1}-x_{2}+x_{3}-x_{4}=0\},
\]
which is $\varphi=-\mathrm{id}$. It is a subuniverse of $\mathbf A^{4}$, being
a subgroup; clause (d) holds because $x_{1}=x_{2}$ iff $x_{3}=x_{4}$;
$\proj_{1,2}(\delta)=\bcol\delta=A^{2}$; and it is symmetric. Yet no bijection
$f\colon\mathbb Z_{3}\to\mathbb Z_{3}$ carries $\delta$ to
$\{x_{1}-x_{2}=x_{3}-x_{4}\}$: such an $f$ would have to satisfy
$f(x_{1})-f(x_{2})=f(x_{3})-f(x_{4})$ whenever $x_{1}-x_{2}+x_{3}-x_{4}=0$, and
putting $(x_{1},x_{2},x_{3},x_{4})=(0,b,b,0)$ gives $2\bigl(f(0)-f(b)\bigr)=0$,
hence $f(0)=f(b)$ for every $b$ since $p$ is odd --- so $f$ is constant, not a
bijection. Since every abelian group of order $3$ is $\mathbb Z_{3}$, no $G$
works, and Lemma~\ref{lem:nice-bridge-abelian} genuinely fails without its
hypothesis.
\end{remark}

\begin{lemma}\label{lem:block-good-bridge}
Let $\delta$ be a bridge from $0_{\mathbf A}$ to $0_{\mathbf A}$ with
$\proj_{1,2}(\delta)\subseteq\bcol\delta$ and $\proj_{1,2}(\delta)$ linked.
Then $\mathbf A$ is BA and centre free.
\end{lemma}

\begin{proof}[Proof outline, with the delicate steps in full]
Replace $\delta$ by $\sigma:=\delta^{\circ}$ of Lemma~\ref{lem:symmetrise} and
put $\omega=\proj_{1,2}(\sigma)=\proj_{1,2}(\delta)$. Note first that
$\bcol\delta$ is \emph{reflexive}: clause (c) of Definition~\ref{def:bridge}
gives $0_{A}\subseteq\proj_{1,2}(\delta)$, and the hypothesis gives
$\proj_{1,2}(\delta)\subseteq\bcol\delta$. Hence $\delta$ is a reflexive
bridge, so Lemma~\ref{lem:symmetrise} applies and gives
$\omega\subseteq\bcol\delta\subseteq\bcol\sigma$ with $\bcol\sigma$
symmetric, which is what legitimises the ``without loss of generality'' below.
Induct on $|A|$.

Suppose $B\subt{T}\mathbf A$ with $T\in\{\TBA,\TC\}$. Since $\omega$ is linked,
$\omega$ meets $(A\setminus B)\times B$ or $B\times(A\setminus B)$; say the
first, and pick $(a,b)$ there.

\emph{Case 1: some subalgebra $D\lneq A$ contains $a$ and $b$.} Put
$\rho=(\omega\circ\omega^{-1})\cap D^{2}$, which is a subalgebra of
$\mathbf D^{2}$, reflexive because $\omega$ is reflexive and symmetric by
construction. Its transitive closure is again a subalgebra, hence a congruence
on $\mathbf D$; let $E$ be the block containing $a$ and $b$, a subuniverse.
Because $\omega$ is reflexive, every undirected $\omega$-edge is a $\rho$-edge
and every $\rho$-edge is a two-edge undirected $\omega$-path, so $\rho$ and the
undirected $\omega$-graph have the same connected components; hence
$\omega\cap E^{2}$ is linked, and $a,b$ lie in one component.
\emph{Note that $(\omega\cup\omega^{-1})\cap D^{2}$ would not serve: a union of
two subuniverses need not be closed.} Moreover
$\proj_{1,2}(\sigma\cap E^{4})=\omega\cap E^{2}$: the inclusion $\subseteq$ is
trivial and $\supseteq$ holds because pair-reflexivity puts
$(x_{1},x_{2},x_{1},x_{2})$ in $\sigma$ for every $(x_{1},x_{2})\in\omega$.
So the inductive hypothesis applies to $\sigma\cap E^{4}$ and makes $E$ BA and
centre free, while Lemma~\ref{lem:ba-center-intersect} makes $B\cap E$ a proper
nonempty subuniverse of $E$ of type $T$ --- proper because $a\in E\setminus B$,
nonempty because $b\in B\cap E$. Contradiction.

\emph{Case 2: no such $D$.} Then $\{a\}\circ\omega=A$, because $\{a\}$ is a
subuniverse by idempotence, hence so is $\{a\}\circ\omega$, and it contains
$a$ and $b$. Put $\xi(x_{1},x_{2})=\sigma(a,x_{1},x_{2},b)$; from
$(a,b,a,b),(a,a,b,b)\in\sigma$ both projections of $\xi$ contain $a$ and $b$,
hence equal $A$. Put $C=B\circ\xi$, which is nonempty because $B$ is and $\proj_{1}(\xi)=A$, so
$C\subte{T}\mathbf A$ by Lemma~\ref{lem:pp-propagation}, with $a\in C$ and
$b\notin C$ --- the latter
because $b\in C$ would give $(a,c,b,b)\in\sigma$ with $c\in B$, and clause (d)
would force $a=c\in B$.

If $T=\TBA$, applying the absorbing term to $(a,b,a,b)$ and $(a,a,b,b)$ gives
$(a,b_{1},b_{2},b)\in\sigma$ with $b_{1},b_{2}\in B$, so $C\cap B\neq
\varnothing$. If $T=\TC$, the three tuples $(a,b,a,b)$, $(a,a,a,a)$ and
$(a,a,b,b)$ witness
\[
  \sigma\cap(\{a\}\times A\times C\times B),\quad
  \sigma\cap(\{a\}\times C\times C\times A),\quad
  \sigma\cap(\{a\}\times C\times A\times B)\;\neq\;\varnothing ,
\]
which is the leave-one-out hypothesis of
Corollary~\ref{cor:no-central-critical} for the ternary relation
$\{(x_{2},x_{3},x_{4}):\sigma(a,x_{2},x_{3},x_{4})\}$ with $n=3$, all types
$\TC$ and all ambients $A$ --- legitimate because $\lll$ is reflexive.
\emph{The three tight boxes of these witnesses --- $\{a\}\times B\times C\times
B$, $\{a\}\times C\times C\times C$, $\{a\}\times C\times B\times B$ --- are
subsets of the displayed ones and do not match the corollary's shape; the
widening is what makes the appeal legitimate.} The corollary's conclusion is
therefore available:
$\sigma\cap(\{a\}\times C\times C\times B)\neq\varnothing$. Put
$C'=\proj_{2}$ of that set, nonempty because the set is. Then
$C'\subte{\TC}A$ by Lemma~\ref{lem:pp-propagation}, using
$\{a\}\circ\omega=A$, and hence
$C'\subte{\TC}C$ by Lemma~\ref{lem:ba-center-intersect}; and by clause (d)
either $a\notin C'$ or $C\cap B\neq\varnothing$.

So in every case we have $\varnothing\neq C'\subt{T}C\subt{T}A$ --- taking
$C'=B\cap C$ when $C\cap B\neq\varnothing$. Then $|C|>1$, and applying the
inductive hypothesis to $\sigma\cap C^{4}$ contradicts $C'\subt{T}C$: its
hypotheses hold because $\{a\}\circ\omega=A$ makes $\omega\cap C^{2}$ a star
centred at $a$, hence linked.
\end{proof}

\begin{longlemma}\label{lem:nontrivial-bridge}
Let $\sigma$ be an irreducible congruence on $\mathbf A$ and $\delta$ a
\emph{pair-reflexive} bridge from $\sigma$ to $\sigma$ with
$\bcol\delta\supsetneq\sigma$. Then
\begin{enumerate}\alphlabels\tightlist
\item $\cover\sigma$ is a congruence;
\item $B/\sigma$ is BA and centre free for every block $B$ of $\cover\sigma$;
\item if $\delta$ is also symmetric then there is a prime $p$ such that every
  block $B$ of $\cover\sigma$ satisfies
  $(B/\sigma;(\delta\cap B^{4})/\sigma)\cong(\mathbb Z_{p}^{n_{B}};
   x_{1}-x_{2}=x_{3}-x_{4})$ for some $n_{B}\ge 0$.
\end{enumerate}
\end{longlemma}

\begin{proof}
By Lemma~\ref{lem:irred-quotient} we may assume $\sigma=0_{\mathbf A}$. Since
$\bcol\delta$ is a subuniverse of $\mathbf A^{2}$ properly containing $\sigma$
and stable under it, minimality gives $\cover\sigma\subseteq\bcol\delta$; the
same argument applied to $\proj_{1,2}(\delta)$, which properly contains $\sigma$
by clause (c), gives $\cover\sigma\subseteq\proj_{1,2}(\delta)$.

Let $B$ be a block of $\LeftLinked(\cover\sigma)$ with $|B|>1$ and
put $\delta'=\delta\cap B^{4}\cap(\cover\sigma\times\cover\sigma)$. Then
$\proj_{1,2}(\delta')=\cover\sigma\cap B^{2}$: for
$(x_{1},x_{2})\in\cover\sigma\cap B^{2}$ pair-reflexivity puts
$(x_{1},x_{2},x_{1},x_{2})$ in $\delta$ --- legitimate because
$\cover\sigma\subseteq\proj_{1,2}(\delta)$ --- and that tuple lies in $B^{4}$
and in $\cover\sigma\times\cover\sigma$ --- \emph{this is where
pair-reflexivity is spent, and without it the projection can collapse to
$0_{B}$}. The resulting relation is linked precisely because $B$ is a block,
and is contained in $\bcol{\delta'}$ because $\cover\sigma\subseteq\bcol\delta$;
so Lemma~\ref{lem:block-good-bridge} makes $B$ BA and centre free. Then
$\cover\sigma\cap B^{2}$, being linked and subdirect on $B$, must be all of
$B^{2}$ by Corollary~\ref{lem:linked-implies-abs}. Blocks of size $1$ satisfy
$B^{2}\subseteq\cover\sigma$ trivially, so
$\cover\sigma=\LeftLinked(\cover\sigma)$ is an equivalence
relation and hence a congruence, giving (a); and (b) follows, one-element
algebras being vacuously BA and centre free.

For (c), fix a block $B$ of $\cover\sigma$ with $|B|\ge2$ and apply
Lemma~\ref{lem:nice-bridge-abelian} to $\delta\cap B^{4}$, whose hypotheses now
hold: $\proj_{1,2}(\delta\cap B^{4})=B^{2}$ by the same pair-reflexivity
computation, and $\bcol{\delta\cap B^{4}}=B^{2}$ because
$B^{2}\subseteq\cover\sigma\subseteq\bcol\delta$. This gives an abelian group
$G_{B}$ with $(B;\delta\cap B^{4})\cong(G_{B};x_{1}-x_{2}=x_{3}-x_{4})$.

It remains to make every $G_{B}$ elementary abelian for one common prime. Put
$\omega=\delta\cap(\cover\sigma\times\cover\sigma)$, which has the same
properties. From $x_{1}-x_{2}=x_{3}-x_{4}$ one pp-defines
$(k+1)x_{1}=k\,x_{2}+x_{3}$ for every $k$, by
\[
  \bigl((k+1)x_{1}=k\,x_{2}+x_{3}\bigr)=\exists x_{4}\;
  \bigl(k\,x_{1}=(k-1)x_{2}+x_{4}\bigr)\wedge(x_{1}-x_{2}=x_{3}-x_{4}),
\]
and hence, identifying $x_{3}$ with $x_{1}$, the relation
$q\,x_{1}=q\,x_{2}$ for every $q$. Let $S$ be what that pp-definition defines
when $x_{1}-x_{2}=x_{3}-x_{4}$ is replaced by $\omega$. If some $G_{B}$ had an
element of order $p_{1}^{2}$, or elements of coprime orders, or if two blocks
had elements of different prime orders, then for a suitable $q$ the relation $S$
would be a subuniverse of $\mathbf A^{2}$ with $S\supsetneq\sigma$ and
$\cover\sigma\not\subseteq S$, contradicting the minimality of $\cover\sigma$.
Hence every element of every $G_{B}$ has one and the same prime order $p$, and
$G_{B}\cong\mathbb Z_{p}^{n_{B}}$.
\end{proof}

\begin{hazard}[It is minimality, not irreducibility]\label{haz:minimality-not-irred}
The contradiction at the end of the proof is with the \emph{minimality of
$\cover\sigma$}, which is a consequence of irreducibility and not the same
statement. The distinction matters wherever $\cover\sigma$ is treated as data
attached to a proof of irreducibility (Caveat~\ref{haz:irreducible}).
\end{hazard}

\begin{proof}[Proof of Lemma~\ref{lem:linear-equiv}]
$(a)\Rightarrow(b)$: take $\delta=S$ of Definition~\ref{def:linear-congruence},
which by Lemma~\ref{lem:linear-bridge} is a bridge with
$\bcol S=\cover\sigma\supsetneq\sigma$.

$(b)\Rightarrow(a)$: given a bridge $\delta$ with
$\bcol\delta\supsetneq\sigma$ --- reflexive by
Lemma~\ref{lem:bridge-reflexive}, since the hypothesis contains
$\bcol\delta\supseteq\sigma$ --- form $\delta^{\circ}$ of
Lemma~\ref{lem:symmetrise}; it is symmetric and pair-reflexive, and
$\bcol{\delta^{\circ}}\supseteq\bcol\delta\supsetneq\sigma$. Now
Lemma~\ref{lem:nontrivial-bridge} gives items (b) and (c) of
Definition~\ref{def:linear-congruence}, with the witness
$S=\delta^{\circ}\cap(\cover\sigma\times\cover\sigma)$.
\end{proof}



\begin{longlemma}\label{lem:pc-inductive-step}
Let $\sigma$ be a congruence on $\mathbf A$ and $\delta$ a reflexive bridge
from $\sigma$ to $\sigma$ with
\begin{enumerate}\romanlabels\tightlist
\item $\delta(x_{1},x_{2},x_{3},x_{4})=\delta(x_{3},x_{4},x_{1},x_{2})$;
\item $(a,b,a,b),(b,a,b,a)\in\delta$ for every $(a,b)\in\proj_{1,2}(\delta)$;
\item $\RightLinked(\proj_{1,2,3}(\delta))=A^{2}$.
\end{enumerate}
Then $(a,a,b,b)\in\delta$ for some $a\neq b$.
\end{longlemma}

\begin{proof}
Factor by $\sigma$ and assume $\sigma=0_{\mathbf A}$; clause (c) of
Definition~\ref{def:bridge} then forces $|A|\ge2$. Write
$P=\proj_{1,2}(\delta)$, which is symmetric by (ii), and
$W=\proj_{1,2,3}(\delta)\sdle\mathbf P\times\mathbf A$. Induct on $|A|$.

\emph{Case 1: some $B\subt{T}\mathbf A$ with $|B|>1$, $T\in\{\TBA,\TC\}$.}
Put $W'=\proj_{1,2,3}(\delta\cap B^{4})$ and write
$W(x_{1},x_{2},x_{3})=\exists x_{4}\,\delta(x_{1},x_{2},x_{3},x_{4})\wedge
A(x_{1})\wedge A(x_{2})\wedge A(x_{3})\wedge A(x_{4})$ with $A$ the full unary
relation. Replacing $A$ by $B$ --- a \emph{single} application of
Lemma~\ref{lem:pp-propagation}, which replaces every occurrence, so no
transitivity of $\subte{\TC}$ is needed --- gives $W'\subte{T}W$. It is
nonempty, containing $(b,b,b)$ for every $b\in B$, and its projections are
$\proj_{1,2}(W')=P\cap B^{2}$, by (ii), and $\proj_{3}(W')=B$. So
Lemma~\ref{lem:absorb-linked-sub} makes $W'$ linked, which is (iii) for
$\delta\cap B^{4}$. Conditions (i) and (ii) are inherited and $\delta\cap B^{4}$
is reflexive. If it is a bridge, the inductive hypothesis applies and produces
the required pair inside $B$. If it is not, then $\proj_{1,2}(\delta\cap
B^{4})=0_{B}$, so by clause (d) every tuple of $\delta\cap B^{4}$ has the form
$(x,x,y,y)$; linkedness of $W'$ on $|B|\ge2$ points then forbids the perfect
matching, so some $(a,a,b,b)\in\delta$ with $a\neq b$.

\emph{Case 2: some $\{a\}\subt{T}\mathbf A$, $T\in\{\TBA,\TC\}$, and no such $B$
of size $>1$.} By (i) the relation $\proj_{1,3}(\delta)$ is symmetric, and by
(iii) it is linked, so $\{a\}\circ\proj_{1,3}(\delta)\neq\{a\}$; pick
$(a,b,c,d)\in\delta$ with $c\neq a$. If $a=b$ then $c=d$ by clause (d) and we
are done. If $\{a\}\circ\proj_{1,3}(\delta)\neq A$ then that set is a strong
subuniverse containing $a$ and $c$, so of size $>1$, which is Case 1.
Otherwise choose $(a,e,b,f)\in\delta$. By (ii), $(b,a,b,a)\in\delta$. Let $g$
be a ternary absorbing operation for $\{a\}$, available for $T=\TC$ by
Lemma~\ref{lem:central-ternary}. Applying $g$
to $(a,e,b,f)$, $(b,a,b,a)$ and $(a,a,a,a)$ gives $(a,a,g(b,b,a),a)\in\delta$,
whence $g(b,b,a)=a$ by clause (d); applying $g$ to $(b,a,b,a)$, $(b,a,b,a)$ and
$(a,e,b,f)$ then gives $(a,a,b,a)\in\delta$, whence $b=a$, a contradiction.

\emph{Case 3: $\mathbf A$ is BA and centre free.} Put
$C=\{(a,b)\in P:\{(a,b)\}\times A\subseteq W\}$. By
Lemma~\ref{lem:bacon-left}, applied to the transpose of $W$,
$C\subte{\TBA}\mathbf P$. If $C=P$ then in particular $(a,a)\in C$ for any $a$;
otherwise $C\subt{\TBA}P$ and Lemma~\ref{lem:absorbing-equality} gives
$C\cap 0_{\mathbf A}\neq\varnothing$. Either way $(a,a)\in C$ for some $a$, so
for every $c$ there is $d$ with $\delta(a,a,c,d)$, and clause (d) forces $d=c$.
Taking $c\neq a$ finishes the proof.
\end{proof}

\begin{lemma}[{\cite[Lem.~7.2]{ZhukJACM}}]\label{lem:absorbing-equality}
If $0_{\mathbf A}\subseteq\sigma\le\mathbf A^{2}$ and
$\omega\subt{\TBA}\sigma$, then $\omega\cap 0_{\mathbf A}\neq\varnothing$.
\end{lemma}

\begin{proof}[Proof of Lemma~\ref{lem:pc-bridges-trivial}]
Put $\xi=\delta\ast\delta^{-1}$, which is symmetric and pair-reflexive. Each of
$\RightLinked(\proj_{1,2,3}(\xi))$ and
$\RightLinked(\proj_{1,2,4}(\xi))$ is a congruence containing
$\sigma$, so by minimality of $\cover\sigma$ each is $\sigma$ or contains
$\cover\sigma$; this gives three cases.

If both equal $\sigma$, then for $(a,b,c,d)\in\xi$ the classes $c/\sigma$ and
$d/\sigma$ are determined by $a/\sigma,b/\sigma$; since $(a,b,a,b)\in\xi$ we get
$(a,c),(b,d)\in\sigma$, so $\xi$ is trivial. \emph{Transferring that
determinacy from $\xi=\delta\ast\delta^{-1}$ to $\delta$ itself is a counting
argument, and it is where the hypothesis
$\proj_{1,2}(\delta)=\proj_{3,4}(\delta)$ is spent:} on
$\cover\sigma/\sigma$, which is finite, $\delta$ induces a bipartite relation
whose left neighbourhoods are nonempty, pairwise disjoint --- that is what the
case gives --- and cover the right side; equal finite cardinalities force every
neighbourhood to be a singleton, so the induced relation is a bijection and
determinacy holds in both directions.

With that, introduce
\begin{align*}
  \zeta_{1}(x_{1},x_{2},x_{3},x_{4})=\;&
    \exists y\exists y'\exists z\exists z'
    \exists t_{1}\exists t_{2}\exists t_{3}\exists t_{4}\;\\
    &\delta(x_{1},y,z,t_{1})\wedge\delta(x_{2},y,z',t_{2})\wedge
     \delta(x_{3},y',z,t_{3})\wedge\delta(x_{4},y',z',t_{4}),\\
  \zeta_{2}(x_{1},x_{2},x_{3},x_{4})=\;&
    \exists y\exists z\;\;\delta(x_{1},y,z,x_{3})\wedge\delta(x_{2},y,z,x_{4}),
\end{align*}
and pick $(a,b,c,d)\in\delta$ with $(a,b)\in\cover\sigma\setminus\sigma$. From
$(a,b,c,d),(b,b,b,b)\in\delta$ one gets $(a,b,a,b)\in\zeta_{1}$ and hence
$\proj_{1,2}(\zeta_{1})\supseteq\cover\sigma$. Split.

\emph{If $\zeta_{1}$ is a bridge}, then $\bcol{\zeta_{1}}=\sigma$ --- forced,
since $\sigma$ is PC --- and $(a,a,c,c)\in\zeta_{1}$ for every
$(a,b,c,d)\in\delta$, so $\proj_{1,3}(\delta)=\sigma$.

\emph{If $\zeta_{1}$ is not a bridge}, there is $(a,a,b,c)\in\zeta_{1}$ with
$(b,c)\notin\sigma$. Take witnesses $y=d$, $y'=d'$, $z=e$, $z'=e'$,
$t_{i}=f_{i}$. Determinacy makes $(e,e')\in\sigma$, so
$(b,d',e,f_{3}),(c,d',e,f_{4})\in\delta$ and therefore
$\proj_{1,2}(\zeta_{2})\supsetneq\sigma$; determinacy in both directions then
makes $\zeta_{2}$ a bridge, whence $\bcol{\zeta_{2}}=\sigma$ and
$\proj_{1,4}(\delta)=\sigma$.

So $\proj_{1,3}(\delta)=\sigma$ or $\proj_{1,4}(\delta)=\sigma$, and by
the same argument with $x_{1},x_{2}$ interchanged, $\proj_{2,4}(\delta)=\sigma$
or $\proj_{2,3}(\delta)=\sigma$. \emph{Two of the four combinations must now be
excluded:} $\proj_{1,3}=\proj_{2,3}=\sigma$ would
give $\proj_{1,2}(\delta)\subseteq\sigma$, against
$\proj_{1,2}(\delta)=\cover\sigma\supsetneq\sigma$, and likewise for
$\proj_{1,4}=\proj_{2,4}$. The two surviving combinations are the conclusion,
the reverse inclusion following from stability under $\sigma$.

If $\RightLinked(\proj_{1,2,3}(\xi))\supseteq\cover\sigma$, choose
a block $B$ of it that is not a $\sigma$-block. Then $a,b,d\in B$ whenever
$c\in B$ and $(a,b,c,d)\in\xi$, so $\xi\cap(A^{2}\times B\times A)\subseteq
B^{4}$ and therefore $\proj_{1,2,3}(\xi\cap B^{4})=\proj_{1,2,3}(\xi)\cap
(A^{2}\times B)$, which is linked. \emph{If $\cover\sigma\cap B^{2}\supsetneq
\sigma\cap B^{2}$ then $\xi\cap B^{4}$ is a bridge and
Lemma~\ref{lem:pc-inductive-step} applies, contradicting PC-ness via
Lemma~\ref{lem:linear-equiv}. If not, the lemma cannot be cited --- it is not a
bridge --- and the case must be inlined: clause (d) then puts both pairs of
every tuple in $\sigma$, and linkedness on a $B$ that is not a $\sigma$-block
produces $(x_{1},x_{1},x_{3},x_{3})\in\xi$ with $(x_{1},x_{3})\notin\sigma$
directly.} The third case is the mirror image, using $\proj_{1,2,4}$.
\end{proof}

\begin{proof}[Proof of Lemma~\ref{lem:no-cross-bridge}]
Let $\sigma_{1},\sigma_{2}$ be congruences on $\mathbf A_{1},\mathbf A_{2}$;
both are irreducible, $\sigma_{1}$ because it is linear
(Definition~\ref{def:linear-congruence}(a)). Replace $\delta$ first by
$\delta\cap(\cover{\sigma_{1}}\times\cover{\sigma_{2}})$, which is again a
bridge with the same collapse, and then by the $\omega$ of
Lemma~\ref{lem:rectangularise}. By clause (d) of that lemma the containment in
$\cover{\sigma_{1}}\times\cover{\sigma_{2}}$ survives, and by clause (c)
neither linking equivalence changes, so the case analysis below is unaffected;
by clause (b), $\bcol\delta$ is now rectangular. Being a bridge inside
$\cover{\sigma_{1}}\times\cover{\sigma_{2}}$, it has
$\proj_{3,4}(\delta)\supsetneq\sigma_{2}$ and stable under $\sigma_{2}$, so
minimality of the cover (Proposition~\ref{prop:cover}) gives
\begin{equation}\label{eq:ncb-proj}
  \proj_{3,4}(\delta)=\cover{\sigma_{2}},
\end{equation}
and symmetrically $\proj_{1,2}(\delta)=\cover{\sigma_{1}}$.

\emph{The easy case.} Suppose $\RightLinked(\bcol\delta)\supsetneq\sigma_{2}$.
The composite $\delta^{-1}\ast\delta$ is a bridge from $\sigma_{2}$ to
$\sigma_{2}$ by Lemma~\ref{lem:bridge-comp} --- both congruences are
irreducible --- with collapse $\bcol\delta^{-1}\circ\bcol\delta$. Iterating,
the $k$-fold composite has collapse $(\bcol\delta^{-1}\circ\bcol\delta)^{k}$,
and for $k$ large this is $\RightLinked(\bcol\delta)\supsetneq\sigma_{2}$. So
Lemma~\ref{lem:linear-equiv} makes $\sigma_{2}$ linear.

So assume $\RightLinked(\bcol\delta)=\sigma_{2}$. Then
$\LeftLinked(\bcol\delta)\supseteq\cover{\sigma_{1}}$: otherwise minimality of
$\cover{\sigma_{1}}$ forces $\LeftLinked(\bcol\delta)=\sigma_{1}$, and
$\bcol\delta$ then induces a bijection
$\mathbf A_{1}/\sigma_{1}\cong\mathbf A_{2}/\sigma_{2}$ carrying
$\cover{\sigma_{1}}/\sigma_{1}$ to $\cover{\sigma_{2}}/\sigma_{2}$, so
$\sigma_{2}$ is linear because $\sigma_{1}$ is.

\emph{Normalising the first and third coordinates.} Put
$\delta'(x_{1},x_{2},x_{3},x_{4})=\delta(x_{1},x_{2},x_{3},x_{4})\wedge
\bcol\delta(x_{1},x_{3})$.

If $\delta'$ is a bridge, it has $\proj_{1,3}(\delta')=\bcol{\delta'}$ by
construction, and we take $\omega=\delta'$.

If $\delta'$ is not a bridge, then --- clauses (a), (b), (d) being inherited ---
clause (c) fails, so there is no $(a,b,c,d)\in\delta$ with
$(a,b)\notin\sigma_{1}$ and $(a,c)\in\bcol\delta$. Put
\[
  \delta''(x_{1},x_{2},x_{3},x_{4})=\exists z\;
    \bcol\delta(x_{1},x_{3})\wedge\delta(x_{2},z,x_{3},x_{4}).
\]
This is a bridge. For clause (d) in one direction: if
$(x_{1},x_{2})\in\sigma_{1}$ then $(x_{2},x_{3})\in\bcol\delta$, so the
failure above applied to $(x_{2},z,x_{3},x_{4})\in\delta$ forces
$(x_{2},z)\in\sigma_{1}$ and hence $(x_{3},x_{4})\in\sigma_{2}$. In the other
direction, if $(x_{3},x_{4})\in\sigma_{2}$ then $(x_{2},z)\in\sigma_{1}$ and
$(x_{2},x_{3})\in\bcol\delta$, whence $(x_{1},x_{2})\in\sigma_{1}$. For
clause (c): since $\delta'\neq\delta$ there is $(a,b,c,d)\in\delta$ with
$(a,c)\notin\bcol\delta$. Some $e$ has $(e,c)\in\bcol\delta$: by
\eqref{eq:ncb-proj} the reflexive pair $(c,c)\in\sigma_{2}\subseteq
\cover{\sigma_{2}}$ is $\proj_{3,4}$ of some $(e,e',c,c)\in\delta$; clause (d)
of Definition~\ref{def:bridge} puts $(e,e')\in\sigma_{1}$, and stability of the
first two variables under $\sigma_{1}$ --- clause (a) --- replaces $e'$ by $e$,
giving $(e,e,c,c)\in\delta$. That $e$ gives $(e,a,c,d)\in\delta''$ with
$(e,a)\notin\sigma_{1}$. Take $\omega=\delta''$.

Either way we have a bridge $\omega$ from $\sigma_{1}$ to $\sigma_{2}$ with
$\proj_{1,3}(\omega)=\bcol\omega=\bcol\delta$.

\emph{The contradiction.} Suppose $\sigma_{2}$ is \emph{not} linear, so it is
PC. Put
\[
  \xi_{1}(x_{1},x_{2},x_{3},x_{4})=\exists x_{5}\exists x_{6}\;
    \omega(x_{5},x_{6},x_{1},x_{2})\wedge\omega(x_{5},x_{6},x_{3},x_{4}),
\]
a bridge from $\sigma_{2}$ to $\sigma_{2}$; by
Lemma~\ref{lem:pc-bridges-trivial} it is
$\sigma_{2}(x_{1},x_{3})\wedge\sigma_{2}(x_{2},x_{4})$ --- the other
alternative is excluded by symmetry of $\xi_{1}$ in its two pairs. Put
\[
  \xi_{2}(x_{1},x_{2},x_{3},x_{4})=\exists x_{5}\exists x_{6}\;
    \omega(x_{5},x_{6},x_{1},x_{2})\wedge\omega(x_{6},x_{5},x_{3},x_{4}).
\]
Here $(x_{5},x_{1}),(x_{6},x_{3})\in\bcol\omega$ and
$(x_{5},x_{6})\in\cover{\sigma_{1}}\subseteq\LeftLinked(\bcol\omega)$, which
gives $(x_{1},x_{3})\in\sigma_{2}$; so $\xi_{2}$ too is
$\sigma_{2}(x_{1},x_{3})\wedge\sigma_{2}(x_{2},x_{4})$ by
Lemma~\ref{lem:pc-bridges-trivial}. Comparing the two,
$(b,a,c,d)\in\omega$ whenever $(a,b,c,d)\in\omega$.

Now put
\[
  \zeta(x_{1},x_{2},x_{3},x_{4})=\exists y_{1}\exists y_{2}\exists y_{3}\;
    \omega(y_{1},y_{2},x_{1},x_{2})\wedge\omega(y_{1},y_{3},x_{1},x_{3})
    \wedge\omega(y_{2},y_{3},x_{1},x_{4}).
\]
If $x_{1}=x_{2}$ then $y_{1}=y_{2}$ by the shape of $\xi_{1}$, and then
$x_{3}=x_{4}$. Take any $(c,d)\in\cover{\sigma_{2}}$; since
$\proj_{3,4}(\omega)=\cover{\sigma_{2}}$ there is
$(a,b)\in\cover{\sigma_{1}}$ with $(a,b,c,d)\in\omega$. Now
$(a,c)\in\bcol\omega$ --- which \emph{is} $(a,a,c,c)\in\omega$, by the
definition of the collapse --- and
$(a,b)\in\cover{\sigma_{1}}\subseteq\LeftLinked(\bcol\omega)$, so
rectangularity in the form of Lemma~\ref{lem:rectangularise}(b) gives
$(b,c)\in\bcol\omega$, that is $(b,b,c,c)\in\omega$. \emph{This is the one
step of the proof that consumes rectangularity.} Sending
$(x_{1},x_{2},x_{3},x_{4})\mapsto(c,d,c,d)$ with $(y_{1},y_{2},y_{3})=(a,b,a)$
witnesses $(c,d,c,d)\in\zeta$; sending it to $(c,c,d,d)$ with
$(y_{1},y_{2},y_{3})=(a,a,b)$ witnesses $(c,c,d,d)\in\zeta$. So
$\zeta(x_{1},x_{2},x_{3},x_{4})\wedge\zeta(x_{3},x_{4},x_{1},x_{2})$ is a
bridge from $\sigma_{2}$ to $\sigma_{2}$ whose collapse contains
$\cover{\sigma_{2}}\supsetneq\sigma_{2}$, so $\sigma_{2}$ is linear by
Lemma~\ref{lem:linear-equiv} --- contradicting the supposition.
\end{proof}

\begin{hazard}[Where rectangularity is spent]\label{haz:rectangularise}
The proof opens by normalising twice, and both steps were asserted in an
earlier draft rather than argued.

\emph{Rectangularity of $\bcol\delta$} is now
Lemma~\ref{lem:rectangularise}: composing $\delta$ with itself enough times
replaces $\bcol\delta$ by $\LeftLinked(\bcol\delta)\circ\bcol\delta$, which is
a union of full boxes, one per connected component, and leaves both linking
equivalences and the containment in
$\cover{\sigma_{1}}\times\cover{\sigma_{2}}$ alone. It is worth knowing that
the normalisation is spent \emph{once}, at $(b,b,c,c)\in\omega$ in the
construction of $\zeta$; nothing else in the proof uses it. Note also that this
is not the situation of Caveat~\ref{haz:reflexivise}, where the property turned
out to hold already: a genuine replacement of $\delta$ is being made here, and
the lemma is what licenses it.

\emph{Choosing $e$ with $(e,c)\in\bcol\delta$} is the second, and it is
three lines from $\proj_{3,4}(\delta)=\cover{\sigma_{2}}$, clause (d) of
Definition~\ref{def:bridge} and stability under $\sigma_{1}$. What made it look
harder than it is was the earlier draft's appeal to $\proj_{1}(\bcol\delta)$,
which is the wrong projection: what is needed is that $c$ lies in
$\proj_{2}(\bcol\delta)$, and the reflexive pair $(c,c)$ supplies it.
\end{hazard}

\begin{proof}[Proof of Lemma~\ref{lem:bridge-to-pc}]
(a) is Lemma~\ref{lem:no-cross-bridge} contrapositively: if $\sigma_{2}$ were
linear then so would $\sigma_{1}$ be, via $\delta^{-1}$.

For (b) and (c), $\delta^{-1}\ast\delta$ is a bridge from $\sigma_{1}$ to
$\sigma_{1}$, and $\sigma_{1}$ is PC, so by
Lemma~\ref{lem:pc-bridges-trivial} it is trivial:
$\sigma_{1}(x_{1},x_{3})\wedge\sigma_{1}(x_{2},x_{4})$ or
$\sigma_{1}(x_{1},x_{4})\wedge\sigma_{1}(x_{2},x_{3})$. Its collapse is
$\bcol\delta\circ\bcol\delta^{-1}$, which is therefore $\sigma_{1}$; so the
relation $\{(a/\sigma_{1},b/\sigma_{2}):(a,b)\in\bcol\delta\}$ has
left-neighbourhoods that are pairwise disjoint. Running the same argument on
$\delta\ast\delta^{-1}$ makes the right-neighbourhoods pairwise disjoint. Both
sides are covered, since $\proj_{1,2}(\delta)=\cover{\sigma_{1}}$ and
$\proj_{3,4}(\delta)=\cover{\sigma_{2}}$, so the induced relation is a
bijection, which is (c), and it is an isomorphism of algebras because
$\bcol\delta$ is a subuniverse --- that is (b).

For (d), put
\[
  \xi(x_{1},x_{2},x_{5},x_{6})=\exists x_{3}\exists x_{4}\;
    \delta(x_{1},x_{2},x_{3},x_{4})\wedge\bcol\delta(x_{5},x_{3})\wedge
    \bcol\delta(x_{6},x_{4}),
\]
a bridge from $\sigma_{1}$ to $\sigma_{1}$ by (c). Lemma~\ref{lem:pc-bridges-trivial}
gives $\xi=\sigma_{1}(x_{1},x_{5})\wedge\sigma_{1}(x_{2},x_{6})$ or
$\sigma_{1}(x_{1},x_{6})\wedge\sigma_{1}(x_{2},x_{5})$, and transporting each
alternative through the bijection of (c) gives the two alternatives of (d).
\end{proof}

\begin{hazard}[The third case]\label{haz:pc-bridges-hedges}
The third case is genuinely the mirror of the second: $(a,b,a,b)\in\xi$ places
$b$ in the block, and $\cover\sigma\subseteq\RightLinked$ places
$a$ and $c$. Writing it out is mechanical, and it is not written out here.
\end{hazard}

\subsection{Types interaction}\label{sec:types-interaction}

\begin{lemma}[{\cite[Lem.~8.19]{ZhukJACM}}]\label{lem:bridge-between-congruences}
Let $\omega,\sigma_{1},\sigma_{2}$ be congruences on $\mathbf A$ with
$\omega\cap\sigma_{1}=\omega\cap\sigma_{2}$ and
$\omega\setminus\sigma_{1}\neq\varnothing$. Then there is a bridge $\delta$
from $\sigma_{1}$ to $\sigma_{2}$ with $\bcol\delta=\sigma_{1}\circ\sigma_{2}$.
\end{lemma}

\begin{proof}
Put
$\delta(x_{1},x_{2},y_{1},y_{2})=\exists z_{1}\exists z_{2}\;
\sigma_{1}(x_{1},z_{1})\wedge\sigma_{1}(x_{2},z_{2})\wedge\omega(z_{1},z_{2})
\wedge\sigma_{2}(z_{1},y_{1})\wedge\sigma_{2}(z_{2},y_{2})$.
Clauses (a) and (b) are visible. For (d): if $(x_{1},x_{2})\in\sigma_{1}$ then
$(z_{1},z_{2})\in\sigma_{1}\cap\omega=\sigma_{2}\cap\omega$, so
$(y_{1},y_{2})\in\sigma_{2}$; \emph{and conversely}, if
$(y_{1},y_{2})\in\sigma_{2}$ then $(z_{1},z_{2})\in\omega\cap\sigma_{2}=
\omega\cap\sigma_{1}$, so $(x_{1},x_{2})\in\sigma_{1}$ --- the direction the
source omits, and the one that spends the other half of the hypothesis. For
(c): $\sigma_{1}\subseteq\proj_{1,2}(\delta)$ by taking $z_{i}=y_{i}=x_{1}$,
and choosing $(a,b)\in\omega\setminus\sigma_{1}$ gives $(a,b,a,b)\in\delta$
with $(a,b)\notin\sigma_{1}$ and, since
$\omega\setminus\sigma_{1}=\omega\setminus\sigma_{2}$, also
$(a,b)\notin\sigma_{2}$. Finally $\bcol\delta=\sigma_{1}\circ\sigma_{2}$: the
inclusion $\supseteq$ by taking $z_{1}=z_{2}$, using reflexivity of $\omega$,
and $\subseteq$ by reading off $(x,z_{1})\in\sigma_{1}$, $(z_{1},y)\in\sigma_{2}$.
\end{proof}

\begin{longlemma}\label{lem:two-stable}
Let $C_{1}\subt[A]{T_{1}(\sigma_{1})}B_{1}\lll A$ and
$C_{2}\subt[A]{T_{2}(\sigma_{2})}B_{2}\lll A$ with
$T_{1},T_{2}\in\{\TBA,\TC,\TS,\TD\}$, $C_{1}\cap B_{2}\neq\varnothing$,
$B_{1}\cap C_{2}\neq\varnothing$ and $C_{1}\cap C_{2}=\varnothing$. Then
$T_{1}=T_{2}\in\{\TBA,\TC,\TD\}$; and if $T_{1}=T_{2}=\TD$ there is a bridge
$\delta$ from $\sigma_{1}$ to $\sigma_{2}$ with
$\bcol\delta=\sigma_{1}\circ\sigma_{2}$.
\end{longlemma}

\begin{proof}
\emph{No $\TS$.} If $T_{1}=\TS$, pick $S_{1}\subseteq C_{1}$ with
$S_{1}\subt{\TBA,\TC}B_{1}$. Since $C_{2}\lll A$ --- the chain for $B_{2}$
extended by one $\TD$ step --- and $B_{1}\cap C_{2}\neq\varnothing$,
Lemma~\ref{lem:intersection-good}(s) gives $S_{1}\cap C_{2}\neq\varnothing$,
hence $C_{1}\cap C_{2}\neq\varnothing$. Symmetrically for $T_{2}$.

\emph{Both absorbing.} If $T_{1},T_{2}\in\{\TBA,\TC\}$ then by
Lemma~\ref{lem:ba-center-intersect} $C_{1}\cap B_{2}\subte{T_{1}}B_{1}\cap
B_{2}$ and $B_{1}\cap C_{2}\subte{T_{2}}B_{1}\cap B_{2}$; both are proper,
since $C_{1}\cap B_{2}=B_{1}\cap B_{2}$ would put $B_{1}\cap C_{2}$ inside
$C_{1}\cap C_{2}=\varnothing$. Their intersection is empty, so
Lemma~\ref{lem:possible-intersections} gives $T_{1}=T_{2}\in\{\TBA,\TC\}$.

\emph{Mixed.} Let $T_{1}\in\{\TBA,\TC\}$ and $T_{2}=\TD$; the mirror case
$T_{1}=\TD$, $T_{2}\in\{\TBA,\TC\}$ is identical with the indices swapped, the
hypotheses being symmetric. By
Lemma~\ref{lem:intersection-good}(d), $(B_{1}\cap B_{2})/\sigma_{2}$ has size
$1$ or is all of $B_{2}/\sigma_{2}$. In the first case the single block is the
one defining $C_{2}$, because $B_{1}\cap C_{2}\neq\varnothing$; then
$B_{1}\cap B_{2}=B_{1}\cap C_{2}$, so $C_{1}\cap B_{2}\subseteq C_{1}\cap
C_{2}=\varnothing$, a contradiction. In the second,
Lemmas~\ref{lem:ba-center-intersect} and \ref{lem:factor-type} give
$(C_{1}\cap B_{2})/\sigma_{2}\subte{T_{1}}B_{2}/\sigma_{2}$, and this is
\emph{proper}: equality would make the $\sigma_{2}$-block defining $C_{2}$ meet
$C_{1}$, i.e.\ $C_{1}\cap C_{2}\neq\varnothing$. So $B_{2}/\sigma_{2}$ has a
proper nonempty strong subuniverse, contradicting
Definition~\ref{def:types}(iii).

\emph{Both $\TD$.} Let $\omega$ be the intersection of the dividing congruences
of the chosen chains for $B_{1}\lll A$ and $B_{2}\lll A$. Apply
Lemma~\ref{lem:intersection-good}(d) to $B_{1}$ and $C_{2}$ with the congruence
$\sigma_{1}$ --- \emph{to $C_{2}$, not to $B_{2}$; against $B_{2}$ the step does
not follow.} Its second alternative, $(B_{1}\cap C_{2})/\sigma_{1}=B_{1}/
\sigma_{1}$, would put a point of $C_{2}$ in the $\sigma_{1}$-block defining
$C_{1}$; so $|(B_{1}\cap C_{2})/\sigma_{1}|=1$ and the ``moreover'' clause
gives $\sigma_{1}\supseteq\omega\cap\sigma_{2}$, because the chain witnessing
$C_{2}\lll A$ is the chain for $B_{2}$ extended by the step
$C_{2}\subt{\TD(\sigma_{2})}B_{2}$ (Caveat~\ref{haz:lll-chain}). Symmetrically
$\sigma_{2}\supseteq\omega\cap\sigma_{1}$, so
$\omega\cap\sigma_{1}=\omega\cap\sigma_{2}$. Finally, for $c_{1}\in C_{1}\cap
B_{2}$ and $c_{2}\in B_{1}\cap C_{2}$ we have $(c_{1},c_{2})\in\omega$, since
both lie in $B_{1}\cap B_{2}$ and each $B_{i}$ sits inside a single block of
every dividing congruence of its own chain; and $(c_{1},c_{2})\notin\sigma_{1}$,
since $c_{2}\in B_{1}$ and $c_{2}\,\sigma_{1}\,c_{1}\in C_{1}$ would give
$c_{2}\in C_{1}\cap C_{2}$. Now Lemma~\ref{lem:bridge-between-congruences}
applies.
\end{proof}

\subsection{Factorisation}\label{sec:factorisation}

\begin{lemma}\label{lem:center-pushed-in}
Let $R\sdle\mathbf A_{1}\times\mathbf A_{2}$ with
$R\cap(B_{1}\times B_{2})\neq\varnothing$, $B_{i}\lll A_{i}$, $\sigma$ a
congruence on $\mathbf A_{1}$ with $B_{1}/\sigma$ BA and centre free. If some
$c\in A_{2}$ has $(E\times\{c\})\cap R\neq\varnothing$ for every $E\in
B_{1}/\sigma$, then such a $c$ may be chosen in $B_{2}$.
\end{lemma}

\begin{proof}[Proof outline]
If $c$ may already be chosen in $B_{2}$ there is nothing to do; otherwise walk
down a chain for $B_{2}\lll A_{2}$ and let $B_{2}''\subt{T(\delta)}B_{2}'$ be
the first step at which the choice fails, so $B_{2}\subseteq B_{2}''$. Put
$S'$ and $S''$ for the sets of $\sigma$-class tuples of length $|A_{1}|$
witnessed by some $c$ in $B_{2}'$, resp.\ $B_{2}''$. A single good $c\in B_{2}'$
makes $S'=(B_{1}/\sigma)^{|A_{1}|}$, and $S''\neq\varnothing$ because
$R\cap(B_{1}\times B_{2})\neq\varnothing$ and $B_{2}\subseteq B_{2}''$ ---
\emph{this is where that hypothesis is spent}. For $T\in\{\TBA,\TC\}$,
Lemmas~\ref{lem:pp-propagation} and \ref{lem:power-to-base} then put a strong
subuniverse on $B_{1}/\sigma$, a contradiction; for $T=\TS$ the same follows
once $R\cap(B_{1}\times D_{2})\neq\varnothing$, which
Lemma~\ref{lem:intersection-good}(s) supplies inside $\mathbf R$.

For $T=\TD$ one builds
$S\le(B_{1}/\sigma)^{|A_{1}|}\times B_{2}'/\delta$ and finds a full column over
$d=c/\delta$, while $(B_{1}/\sigma)^{|A_{1}|}\times B_{2}''/\delta\not\subseteq
S$ by minimality. Corollary~\ref{cor:intersection-good} then computes
$\proj_{|A_{1}|+1}(S)$. It offers three alternatives and only the third gives
the contradiction. ``Empty'' is excluded by
$R\cap(B_{1}\times B_{2})\neq\varnothing$. ``Size one'' is excluded thus: the
single $\delta$-class would have to be $[c]_{\delta}$, since $c$ has an
$R$-partner in $B_{1}$; and $B_{2}''=B_{2}'\cap E$ for a $\delta$-block $E$, so
either $E=[c]_{\delta}$, putting $c$ in $B_{2}''$ against the minimal choice of
$B_{2}'$, or $E\cap[c]_{\delta}=\varnothing$, making
$R\cap(B_{1}\times B_{2}'')=\varnothing$ against $B_{2}\subseteq B_{2}''$. So
$S$ is a central relation, and Lemmas~\ref{lem:central-relation} and
\ref{lem:power-to-base} contradict the hypotheses.
\end{proof}

\begin{lemma}\label{lem:leftlinked-stay-full}
Under the hypotheses of Lemma~\ref{lem:center-pushed-in}, if
$(\LeftLinked(R)\cap B_{1}^{2})/\sigma=(B_{1}/\sigma)^{2}$ then
$\LeftLinked\bigl(R\cap(B_{1}\times B_{2})\bigr)/\sigma=
(B_{1}/\sigma)^{2}$.
\end{lemma}

\begin{proof}[Proof outline]
Write $R_{n}=(R\circ R^{-1})^{n}$, which is reflexive and symmetric because $R$
is subdirect, so $R_{2n}=R_{n}\circ R_{n}^{-1}$ --- \emph{this is why the
source restricts to $n=2^{k}$}. If $(R_{1}\cap B_{1}^{2})/\sigma$ is already
full, Lemma~\ref{lem:bacon-left} applied to
$S=\{(a/\sigma,b):a\in B_{1},(a,b)\in R\}$, subdirect over
$(B_{1}/\sigma)\times\proj_{2}(S)$ rather than over $(B_{1}/\sigma)\times A_{2}$,
produces the required $c$, which Lemma~\ref{lem:center-pushed-in} moves into
$B_{2}$; two classes are then linked through $c$. Otherwise take the largest
$n=2^{k}$ at which fullness fails, apply Lemma~\ref{lem:bacon-left} to $R_{n}$,
push $c$ into $B_{1}$, and conclude by Lemma~\ref{lem:central-relation} that
$(R_{n}\cap B_{1}^{2})/\sigma$ is full after all.
\end{proof}

\begin{longlemma}\label{lem:cut-or-nothing}
Let $\sigma$ be a dividing congruence for $B\lll A$, let $\delta$ be a
congruence on $\mathbf A$ with $(\delta\cap B^{2})/\sigma\neq B^{2}/\sigma$, and
let $\omega$ be the intersection of the dividing congruences of the chosen chain
for $B\lll A$. Then \textup{(1)} $\delta\cap B^{2}\subseteq\sigma\cap B^{2}$;
\textup{(2)} $\sigma\supseteq\delta\cap\omega$; \textup{(3)}
$(\delta\vee(\sigma\cap\omega))\cap B^{2}=\sigma\cap B^{2}$; \textup{(4)}
$(\delta\vee(\sigma\cap\omega))\cap\omega=\sigma\cap\omega$.
\end{longlemma}

\begin{proof}
Throughout we use $B^{2}\subseteq\omega$, which holds because each $B_{i}$ in
the chain lies inside a single block of its own dividing congruence, and $B$ is
contained in every $B_{i}$. \emph{It is used at least four times below.}

(2). The hypothesis gives $\sigma$-classes $C_{1},C_{2}$ meeting $B$ with
$\bigl((C_{1}\cap B)\times(C_{2}\cap B)\bigr)\cap\delta=\varnothing$. Then
$C_{1}\cap B\subt{\TD(\sigma)}B$, and Lemma~\ref{lem:self-intersection-pc} with
$m=k$ and $B_{k}=C_{1}\cap B$ gives
$\bigl((C_{1}\cap B)\circ\delta\cap B\bigr)/\sigma=\{C_{1}\}$, since the value
$B/\sigma$ is excluded by the choice of $C_{2}$; its second clause is
$\sigma\supseteq\delta\cap\omega$.

(1) follows, since $\delta\cap B^{2}\subseteq\delta\cap\omega\subseteq\sigma$.

(4). Put $\delta'=\delta\vee(\sigma\cap\omega)$, which is
$\LeftLinked(\delta\circ(\sigma\cap\omega))$. If
$(\delta'\cap B^{2})/\sigma$ were full then Lemma~\ref{lem:leftlinked-stay-full}
would make $\LeftLinked$ of the restricted relation full, so some
single linking step would cross $\sigma$-classes: there would be
$a_{1},a_{2},c\in B$ and $b_{1},b_{2}\in A$ with $(a_{i},b_{i})\in\delta$,
$(b_{i},c)\in\sigma\cap\omega$ and $(a_{1},a_{2})\notin\sigma$. But
$a_{1},a_{2},c\in B$ gives $(a_{i},c)\in\omega$, hence $(a_{i},b_{i})\in\omega$,
hence $(a_{i},b_{i})\in\sigma$ by (2), hence $(a_{1},a_{2})\in\sigma$ --- a
contradiction. So $(\delta'\cap B^{2})/\sigma\neq(B/\sigma)^{2}$, and applying
(2) to $\delta'$ gives $\sigma\supseteq\delta'\cap\omega$; with
$\delta'\supseteq\sigma\cap\omega$ this is (4).

(3) follows from (4) by intersecting with $B^{2}\subseteq\omega$.
\end{proof}

\begin{lemma}\label{lem:main-existence}
Let $B\lll A$ and let $\sigma$ be a congruence on $\mathbf A$ with
$|B/\sigma|>1$ and $B/\sigma$ BA and centre free. Then there is a dividing
congruence $\delta\supseteq\sigma$ for $B\le A$.
\end{lemma}

\begin{proof}
Choose $\delta\supseteq\sigma$ maximal with $|B/\delta|>1$; it exists because
$\sigma$ qualifies. Then $B/\delta$ is BA and centre free, since the preimage
under $B/\sigma\twoheadrightarrow B/\delta$ of a proper nonempty strong
subuniverse would be one of $B/\sigma$.

It remains to see that $\delta$ is irreducible. Suppose
$\delta=S_{1}\cap\dots\cap S_{k}$ with each $S_{i}\supsetneq\delta$ a
subuniverse of $\mathbf A^{2}$ \emph{stable under $\delta$}. Some $S_{i}$ has
$B^{2}\not\subseteq S_{i}$, else $B^{2}\subseteq\delta$. Now
$\LeftLinked(S_{i})$ is a congruence containing
$S_{i}\supsetneq\delta\supseteq\sigma$, so maximality of $\delta$ forces
$B^{2}\subseteq\LeftLinked(S_{i})$, i.e.\
$(\LeftLinked(S_{i})\cap B^{2})/\delta=(B/\delta)^{2}$. By
Lemma~\ref{lem:leftlinked-stay-full}, applied \emph{with the congruence
$\delta$}, $\LeftLinked(S_{i}\cap B^{2})/\delta=(B/\delta)^{2}$,
so $(S_{i}\cap B^{2})/\delta$ is linked. It is also \emph{proper}: stability
under $\delta$ makes $S_{i}$ a union of products of $\delta$-blocks, so
$(S_{i}\cap B^{2})/\delta=(B/\delta)^{2}$ would give $B^{2}\subseteq S_{i}$.
Corollary~\ref{lem:linked-implies-abs} then produces a proper nonempty strong
subuniverse of $B/\delta$, a contradiction.
\end{proof}

\begin{hazard}[Where each hypothesis is spent]\label{haz:main-existence}
Two things in that proof are easy to lose. The relation must be factored by
$\delta$ and not by $\sigma$ --- $(S_{i}\cap B^{2})/\sigma$ lives on $B/\sigma$
and the conclusion is about $B/\delta$, and it is $B/\delta$ that was shown BA
and centre free. And properness needs the words \emph{stable under $\delta$} in
the unfolding of irreducibility: that clause is what makes $S_{i}$ a union of
products of $\delta$-blocks, and it is the only part of
Definition~\ref{def:irreducible} doing work here. Without it,
$S_{i}\cap B^{2}\subsetneq B^{2}$ survives factoring only by accident.
\end{hazard}

\begin{lemma}\label{lem:factor-by-delta}
Let $\delta$ be a congruence on $\mathbf A$ and
$C\subt[A]{T(\sigma)}B\lll A$ with $T\in\{\TPC,\TL\}$. Then
$C/\delta=B/\delta$, or $C/\delta\subt{\TS}B/\delta$, or
$C/\delta\subt[A/\delta]{T}B/\delta$; and in the last case
$\delta\cap B^{2}\subseteq\sigma\cap B^{2}$.
\end{lemma}

\begin{proof}[Proof outline]
Put $R=\{(a/\sigma,a/\delta):a\in B\}$ and note
$R\circ R^{-1}=(\delta\cap B^{2})/\sigma$, so
$\LeftLinked(R)=(B/\sigma)^{2}$ exactly when
$(\delta\cap B^{2})/\sigma=B^{2}/\sigma$. In that case
Lemma~\ref{lem:bacon-left} yields $U\subte{\TBA,\TC}B/\delta$ with
$(B/\sigma)\times U\subseteq R$, and $U\subseteq C/\delta$, giving the first two
alternatives.

Otherwise Lemma~\ref{lem:cut-or-nothing}(4) gives
$\sigma\cap\omega=\delta'\cap\omega$ for $\delta'=\delta\vee(\sigma\cap\omega)$,
and $B/\delta'\cong B/\sigma$ by (3), so $B/\delta'$ is BA and centre free and
Lemma~\ref{lem:main-existence} produces a dividing congruence
$\delta''\supseteq\delta'$ for $B$. Applying
Lemma~\ref{lem:cut-or-nothing}(2) to $\delta''$ --- legitimate because
$\sigma\cap B^{2}\subseteq\delta''$ and $|B/\delta''|>1$ give
$(\delta''\cap B^{2})/\sigma\neq B^{2}/\sigma$ --- yields
$\delta''\cap\omega=\sigma\cap\omega$, hence $\delta''\cap B^{2}=\sigma\cap
B^{2}$, so the $\delta''$-classes and $\sigma$-classes agree on $B$ and
$C/\delta\subt[A/\delta]{\mathcal T(\delta''/\delta)}B/\delta$. Finally
$\mathcal T=T$: $\omega\setminus\delta''\neq\varnothing$, since
$\omega\supseteq B^{2}$ and $|B/\sigma|>1$, so
Lemma~\ref{lem:bridge-between-congruences} gives a bridge from $\delta''$ to
$\sigma$, and Lemma~\ref{lem:no-cross-bridge} in both directions makes the two
congruences of the same type. The last clause holds because
$\delta\cap B^{2}\subseteq\delta'\cap B^{2}=\sigma\cap B^{2}$.
\end{proof}

\subsection{The headline properties}\label{sec:headline}

\begin{proof}[Proof of Lemma~\ref{lem:ubiquity}]
If $\mathbf B$ has a proper nonempty binary absorbing or central subuniverse,
take it. Otherwise $\mathbf B$ is BA and centre free and $|B|>1$, so
Lemma~\ref{lem:main-existence} with $\sigma=0_{\mathbf A}$ gives a dividing
congruence $\delta$ for $B$; take $C=B\cap E$ for any block $E$ \emph{of
$\delta$} with $B\cap E\neq B$, which exists because $|B/\delta|>1$. Its type is
$\TL$ or $\TPC$ according to $\delta$.
\end{proof}

\begin{proof}[Proof of Lemma~\ref{lem:intersect-all} and
  Corollary~\ref{cor:intersect-all-weak}]
(t). For $T\in\{\TBA,\TC\}$ this is Lemma~\ref{lem:ba-center-intersect}. For
$T=\TS$, take $E\subseteq C$ with $E\subt{\TBA,\TC}B$; if $C\cap D=\varnothing$
there is nothing to prove, and otherwise $B\cap D\neq\varnothing$, so
Lemma~\ref{lem:intersection-good}(s) gives $E\cap D\neq\varnothing$ and
Lemma~\ref{lem:ba-center-intersect} gives $E\cap D\subt{\TBA,\TC}B\cap D$. For
$T\in\{\TPC,\TL\}$, Lemma~\ref{lem:intersection-good}(d) makes $(B\cap D)/\sigma$
empty, a point, or all of $B/\sigma$; the three cases give $C\cap D=\varnothing$,
$C\cap D\in\{\varnothing,B\cap D\}$, and $C\cap D\subt{T(\sigma)}B\cap D$
respectively.

(i). Apply (t) to each step of a chain for $B\lll^{A}A$, obtaining
$B\cap D\dlll^{A}D$. The corollary follows by composing with $D\lll A$.
\end{proof}

\begin{hazard}[The superscript in (i) is needed]\label{haz:intersect-all-strength}
Corollary~\ref{cor:intersect-all-weak} is weaker than
Lemma~\ref{lem:intersect-all}(i), and three later arguments need the stronger
form: Corollary~\ref{cor:prop-relations}(b) and (m), which conclude
$R''\dlll^{R}R'$, and Lemma~\ref{lem:multitype-chain}. Stating only the
corollary --- as is natural, since it is the symmetric-looking form --- makes
all three unprovable.
\end{hazard}

\begin{proof}[Proof of Lemma~\ref{lem:reverse-hom}]
(bt) is proved directly, from the definitions; the reflection corollary of
\S\ref{sec:properties} is \emph{derived} from this lemma and so may not be used
here. For
$T\in\{\TBA,\TC\}$ the witnessing term transfers: $f$ is a homomorphism, so
$f\bigl(t(\bar a)\bigr)=t\bigl(f(\bar a)\bigr)$, and a tuple with all but one
entry in $f^{-1}(C')$ and the remaining entry in $f^{-1}(B')$ maps to one with
all but one entry in $C'$ and the remaining entry in $B'$; the centrality
clause pulls back the same way, since $f\bigl(\Sg(X)\bigr)=\Sg\bigl(f(X)\bigr)$
for surjective $f$. For $T=\TS$ the witness $D'\subseteq C'$ pulls back to
$f^{-1}(D')\subseteq f^{-1}(C')$ and the previous case applies. For
$T\in\{\TD,\TL,\TPC\}$, $C'=B'\cap E$ for a block $E$ of $\sigma$, so
$f^{-1}(C')=f^{-1}(B')\cap f^{-1}(E)$ with $f^{-1}(E)$ a block of
$f^{-1}(\sigma)$; condition (i) of Definition~\ref{def:types} transfers because
$f^{-1}\bigl(\cover\sigma\bigr)=\cover{f^{-1}(\sigma)}$ by
Lemma~\ref{lem:irred-quotient}, and condition (iii) because
$f^{-1}(B')/f^{-1}(\sigma)\cong B'/\sigma$.

(b) follows by iterating (bt) along a chain, and (bm) from (bt) by
intersecting.
\end{proof}

\begin{proof}[Proof of Lemma~\ref{lem:propagation}]
The backward clauses are Lemma~\ref{lem:reverse-hom}.

(fs) is Lemma~\ref{lem:factor-type}. (ft) splits by type: for
$T\in\{\TBA,\TC,\TS\}$ it follows from (fs); for $T\in\{\TPC,\TL\}$ it is
Lemma~\ref{lem:factor-by-delta}; and for $T=\TD$, which that lemma does not
cover, split the witnessing congruence into its linear and PC cases and apply
the lemma to each, the two conclusions being the same three alternatives. Then
(f) follows by iteration, all three outcomes of (ft) staying inside the type
list of Definition~\ref{def:lll}.

(fm). Let $C=C_{1}\cap\dots\cap C_{t}$ with $C_{i}\subt[A]{T(\sigma_{i})}B$ and
induct on $t$. Write $\delta$ for the kernel of $f$, so that $f(X)$ may be read
as $X/\delta$ throughout. For $t=1$ this is (ft), and \emph{the hypothesis that
$f(B)$ is $\TS$-free is exactly what kills its middle alternative}, which
$\MT T$ cannot absorb. If $C_{i}/\delta=B/\delta$ for some $i$, say $i=t$, put
$D_{j}=C_{j}\cap C_{t}$; by Lemma~\ref{lem:intersect-all}(t)
$D_{j}\subte[A]{T(\sigma_{j})}C_{t}$, and the inductive hypothesis applied to
$C\subte{\MT T}C_{t}$ gives the claim. Otherwise every
$C_{i}/\delta\subt[A/\delta]{T}B/\delta$, and
$\delta\cap B^{2}\subseteq\sigma_{i}\cap B^{2}$ for every $i$ by the last clause
of Lemma~\ref{lem:factor-by-delta}. That is what makes the images meet
correctly: $(C_{i}\circ\delta)\cap B=C_{i}$, so a $\delta$-class $E$ lying in
every $C_{i}/\delta$ has $E\cap B\subseteq C_{1}\cap\dots\cap C_{t}$, whence
\[
  (C_{1}\cap\dots\cap C_{t})/\delta=C_{1}/\delta\cap\dots\cap C_{t}/\delta ,
\]
the inclusion $\subseteq$ being trivial. This is the one place where the failure
of images to commute with intersections (\S\ref{sec:legislation}, item
\ref{leg:quotient}) is repaired, and the saturation hypothesis
$\delta\cap B^{2}\subseteq\sigma_{i}\cap B^{2}$ is precisely what buys it.
\end{proof}

\begin{longtheorem}[Stable intersection, restated]
Let $C_{i}\subt[A]{T_{i}(\sigma_{i})}B_{i}\lll A$ with
$T_{i}\in\{\TBA,\TC,\TS,\TL,\TPC\}$ for $i\in[n]$, $n\ge2$, and suppose
$\bigcap_{i}C_{i}=\varnothing$ while $B_{j}\cap\bigcap_{i\neq j}C_{i}\neq
\varnothing$ for every $j$. Then all $T_{i}$ are $\TBA$; or all are $\TL$ and
for all $k,\ell$ there is a bridge $\delta$ from $\sigma_{k}$ to $\sigma_{\ell}$
with $\bcol\delta=\sigma_{k}\circ\sigma_{\ell}$; or $n=2$ and
$T_{1}=T_{2}=\TC$; or $n=2$, $T_{1}=T_{2}=\TPC$ and $\sigma_{1}=\sigma_{2}$.
\end{longtheorem}

\begin{proof}[Proof of Theorem~\ref{thm:stable-intersection}]
For $n=2$ combine Lemmas~\ref{lem:two-stable} and \ref{lem:no-cross-bridge}.
For general $n$ put $B_{2}'=B_{2}\cap C_{3}\cap\dots\cap C_{n}$ and
$C_{2}'=C_{2}\cap B_{2}'$; then $C_{2}'\subt{T_{2}(\sigma_{2})}B_{2}'\lll A$ by
Lemma~\ref{lem:intersect-all}, properly because $C_{2}'=B_{2}'$ would make
$C_{1}\cap B_{2}'$ empty. The three hypotheses of the pair case hold, so
$T_{1}=T_{2}$ and, for the dividing types, the bridge exists.

If $T_{1}=\TPC$ then $\sigma_{1}\circ\sigma_{2}=\sigma_{1}=\sigma_{2}$. Indeed,
composing the bridge with its transpose gives a bridge from $\sigma_{1}$ to
itself, and if $\sigma_{1}\circ\sigma_{2}\supsetneq\sigma_{1}$ its collapse
properly contains $\sigma_{1}$, making $\sigma_{1}$ linear by
Lemma~\ref{lem:linear-equiv}. That leaves the possibility of proper nesting,
$\sigma_{1}\circ\sigma_{2}=\sigma_{1}$ with $\sigma_{1}\neq\sigma_{2}$, which
is excluded by Lemma~\ref{lem:bridge-to-pc}: a bridge between PC congruences
makes $\mathbf A/\sigma_{1}\cong\mathbf A/\sigma_{2}$ with the induced relation
on blocks bijective, and two congruences on one finite algebra with equal
quotient sizes and $\sigma_{2}\subseteq\sigma_{1}$ are equal. With $\sigma_{1}=\sigma_{2}=:\sigma$, the blocks $E_{1},E_{2}$ defining
$C_{1},C_{2}$ are \emph{distinct} --- equal blocks would put the point of
$B_{1}\cap C_{2}'$ in $C_{1}\cap C_{2}'$ --- so $C_{1}\cap C_{2}=\varnothing$;
applying this to every pair and then the hypothesis at $j=3$ forces $n=2$.

For $T_{1}=\TC$ the argument just given has no analogue --- there is no
congruence and no pair of blocks --- and $n=2$ is instead
Lemma~\ref{lem:no-central-critical}, whose two hypotheses are exactly (ii) and
(iii) here.
\end{proof}

\begin{hazard}[What the type-$\TC$ case costs]\label{haz:d4}
The type-$\TC$ argument is the one place in this section that reaches outside
the material developed here, to \cite[Lem.~6.11]{ZhukStrong}: \emph{for
$n\ge3$ and $E_{i}$ central in $\mathbf D$, there is no
$(E_{1},\dots,E_{n})$-essential relation $R\le D^{n}$}. Two notes for anyone
checking it. Its printed proof cites \emph{itself} three times where it means
the preceding lemma, so a self-contained rendering must supply both. And its
ingredients are exactly the Barto--Kazda criterion
(Lemma~\ref{lem:barto-kazda}) and ``central implies ternary absorbing''
(Lemma~\ref{lem:central-ternary}), both of which are already available above.

Nonemptiness and properness of each $E_{i}$ are established before the
essential-relation theorem is applied, and both are needed: essentiality is a
statement about proper nonempty subsets.
\end{hazard}

\begin{proof}[Proof of Corollary~\ref{cor:stable-intersection}]
Pull back along the projections $f_{i}:\mathbf R\to\mathbf A_{i}$ using
Lemma~\ref{lem:reverse-hom}(b),(bt), apply
Theorem~\ref{thm:stable-intersection} inside $\mathbf R$, and translate
$\bcol\delta=\sigma_{k}'\circ\sigma_{\ell}'$ back as
$\sigma_{k}\circ\proj_{k,\ell}(R)\circ\sigma_{\ell}$.
\end{proof}

\begin{proof}[Proof of Lemma~\ref{lem:multitype-chain}]
With $C=\bigcap_{i\le n}C_{i}$ and $D_{j}=\bigcap_{i\le j}C_{i}$,
Lemma~\ref{lem:intersect-all}(t) gives $D_{j+1}\subte[A]{T}D_{j}$, using
$D_{j}\lll A$ from part (i); collapsing the equalities gives the chain.
\end{proof}

\begin{proof}[Proof of Lemma~\ref{lem:anchored-intersection}]
By Lemma~\ref{lem:multitype-chain} fix a chain
$C_{2}=F_{s}\subt[A_{2}]{T}\dots\subt[A_{2}]{T}F_{0}=B_{2}$. If $s=0$ then
$C_{2}=B_{2}$ and the second part of (i) is the conclusion. Each truncation
of the chain witnesses $F_{j}\lll^{A_{2}}B_{2}$, and concatenating it with
the chain for $B_{2}\lll A_{2}$ gives $F_{j}\lll A_{2}$ for every $j$.

Suppose for contradiction that $W\cap(C_{1}\times C_{2})=\varnothing$. Along
the chain, $W\cap(C_{1}\times F_{0})\neq\varnothing$ by (i) and
$W\cap(C_{1}\times F_{s})=\varnothing$, so there is a least $j$ with
$W\cap(C_{1}\times F_{j})=\varnothing$; fix it, so that
$W\cap(C_{1}\times F_{j-1})\neq\varnothing$. Write the step as
$F_{j}\subt[A_{2}]{T(\xi)}F_{j-1}$ and let $T'$ be the class of $\xi$ ---
linear or PC, exhaustively, by Definition~\ref{def:linear-congruence} --- so
that $F_{j}\subt[A_{2}]{T'(\xi)}F_{j-1}$; for $T\in\{\TL,\TPC\}$ the class
is $T'=T$, and for $T=\TD$ it is whichever the step's congruence has.

Apply Corollary~\ref{cor:stable-intersection} to $W$ with the two-member
family
\[
  C_{1}\subt[A_{1}]{T_{1}(\sigma_{1})}B_{1}\lll A_{1},
  \qquad
  F_{j}\subt[A_{2}]{T'(\xi)}F_{j-1}\lll A_{2}.
\]
Its hypotheses hold: $W$ is subdirect; the joint intersection
$W\cap(C_{1}\times F_{j})$ is empty by the choice of $j$; and both
leave-one-out configurations are nonempty ---
$W\cap(B_{1}\times F_{j})\supseteq W\cap(B_{1}\times C_{2})\neq\varnothing$
by (i), since $C_{2}=F_{s}\subseteq F_{j}$, and
$W\cap(C_{1}\times F_{j-1})\neq\varnothing$ by the minimality of $j$. So one
of the four alternatives holds of the pair $\{T_{1},T'\}$. None can:
$T'\in\{\TL,\TPC\}$, so (ba) and (c) fail; $\TS$ occurs in no alternative;
and (l), (pc) require $T_{1}=T'\in\{\TL,\TPC\}$, which (ii) excludes ---
for $T_{1}=\TL$ it puts $T=\TPC$, making every step of the chain of class
$\TPC$, and symmetrically for $T_{1}=\TPC$. The contradiction refutes the
assumption, so $W\cap(C_{1}\times C_{2})\neq\varnothing$.
\end{proof}

\begin{proof}[Proof of Lemma~\ref{lem:linear-on-top}]
By Lemma~\ref{lem:linear-equiv} there is a bridge $\delta$ from $\sigma$ to
$\sigma$ with $\bcol\delta\supsetneq\sigma$, reflexive by
Lemma~\ref{lem:bridge-reflexive}; replace it by $\delta^{\circ}$, symmetric and
pair-reflexive.
Lemma~\ref{lem:nontrivial-bridge}(c) applies, and $\cover\sigma=A^{2}$ is a
congruence with the single block $A$, so $A/\sigma\cong\mathbb Z_{p}^{m}$ as a
group. Now $m=0$ gives $\sigma=A^{2}$, impossible for an irreducible congruence
(\S\ref{sec:legislation}, item \ref{leg:emptyfamily}). And $m\ge2$ contradicts
irreducibility: $A/\sigma$ is affine, so the basic operation acts as
$\sum_{i}c_{i}x_{i}$ with $\sum_{i}c_{i}=1$, and the coordinate congruences
$\theta_{i}=\{(x,y):x_{i}=y_{i}\}$ are linear subspaces of
$(\mathbb Z_{p}^{m})^{2}$, hence subuniverses of $(\mathbf A/\sigma)^{2}$, each
properly containing $0_{\mathbf A/\sigma}$, with
$\bigcap_{i}\theta_{i}=0_{\mathbf A/\sigma}$ --- so
$0_{\mathbf A/\sigma}$ is reducible, and by Lemma~\ref{lem:irred-quotient} so is
$\sigma$. Hence $m=1$, and the group isomorphism is an isomorphism of algebras
by Proposition~\ref{prop:p-divides}.
\end{proof}

\begin{lemma}\label{lem:multiply-all-linear}
Let $C_{1}\subt[A]{\MT T}B_{1}\lll A$ with $T\in\{\TPC,\TL,\TD\}$, let
$B_{2}\lll A$, $C_{1}\cap B_{2}=\varnothing$ and $B_{1}\cap B_{2}\neq
\varnothing$. Then $\bigl(C_{1}\circ(\omega_{1}\cap\dots\cap\omega_{s})\bigr)
\cap B_{2}=\varnothing$, where $\omega_{1},\dots,\omega_{s}$ are all congruences
of type $T$ on $\mathbf A$ with $\cover{\omega_{i}}\supseteq B_{1}^{2}$.
\end{lemma}

\begin{proof}[Proof of Case~1, and of Case~2 from Claim~\ref{cl:multiply-step}]
The induction is downward on $|B_{1}|$: the inductive hypothesis is the lemma
for every larger first argument, and it starts at $B_{1}=A$.

By Definition~\ref{def:multitypes} write $C_{1}=\bigcap_{i\in[t]}C_{1}^{i}$
with $t\ge1$ and $C_{1}^{i}\subt[A]{T(\sigma_{i})}B_{1}$ --- of type $T$, not
merely of type $\TD$, which is what makes each $\sigma_{i}$ a congruence of type
$T$ with $\cover{\sigma_{i}}\supseteq B_{1}^{2}$ and hence one of the
$\omega_{j}$. Put $K_{i}=C_{1}^{i}\circ\sigma_{i}$ and
$\Omega=\omega_{1}\cap\dots\cap\omega_{s}$. Definition~\ref{def:types}(ii)
gives $C_{1}^{i}=K_{i}\cap B_{1}$, and
Corollary~\ref{cor:prop-times-cong}(e),(f) gives
$K_{i}\subt[A]{T(\sigma_{i})}B_{1}\circ\sigma_{i}\lll A$, so in particular
$K_{i}\lll A$. Two inclusions are used, and both are immediate:
\begin{equation}\label{eq:mal-incl}
  C_{1}\circ\Omega\;\subseteq\;\bigcap_{i\in[t]}K_{i},
  \qquad\text{and}\qquad
  X\subseteq Y\;\Longrightarrow\;X\circ\Omega\subseteq Y\circ\Omega .
\end{equation}
The first holds because $C_{1}\subseteq C_{1}^{i}$ and
$\Omega\subseteq\sigma_{i}$ for each $i$, the latter because $\sigma_{i}$ is
one of the $\omega_{j}$.

Put $D=\bigcap_{i\in[t]}K_{i}\cap B_{2}$, and note
$D\cap B_{1}=\bigcap_{i}(K_{i}\cap B_{1})\cap B_{2}=C_{1}\cap B_{2}=
\varnothing$.

\emph{Case 1: $D=\varnothing$.} Then by \eqref{eq:mal-incl}
$(C_{1}\circ\Omega)\cap B_{2}\subseteq\bigcap_{i}K_{i}\cap B_{2}=\varnothing$,
which is the conclusion. This is the whole of the case $B_{1}=A$, since there
$D=D\cap B_{1}=\varnothing$.

\emph{Case 2: $D\neq\varnothing$.} Each $K_{i}\lll A$ and $B_{2}\lll A$, so
Corollary~\ref{cor:intersect-all-weak}, applied $t$ times, gives
$D\dlll A$; as $D\neq\varnothing$, $D\lll A$. Fix the chain
$A=X_{0}\supsetneq X_{1}\supsetneq\dots\supsetneq X_{r}=B_{1}$ witnessing
$B_{1}\lll A$. Then $D\cap X_{0}=D\neq\varnothing$ and $D\cap X_{r}=
D\cap B_{1}=\varnothing$, so there is a least $p\in[r]$ with
$D\cap X_{p}=\varnothing$. Write $B_{1}'=X_{p}$ and $B_{1}''=X_{p-1}$, so that
\begin{equation}\label{eq:mal-step}
  B_{1}\subseteq B_{1}'\subt[A]{}B_{1}''\lll A,\qquad
  D\cap B_{1}'=\varnothing,\qquad D\cap B_{1}''\neq\varnothing .
\end{equation}

\begin{claim}\label{cl:multiply-step}
That step is of multi-type $T$: $B_{1}'\subt[A]{\MT T}B_{1}''$.
\end{claim}

Granting the claim, the inductive hypothesis applies to
$B_{1}'\subt[A]{\MT T}B_{1}''\lll A$ and $D\lll A$, whose hypotheses
$B_{1}'\cap D=\varnothing$ and $B_{1}''\cap D\neq\varnothing$ are
\eqref{eq:mal-step}, and which sits at a strictly larger first argument since
$B_{1}\subseteq B_{1}'\subsetneq B_{1}''$. It returns
$(B_{1}'\circ\Omega'')\cap D=\varnothing$, where $\Omega''$ is the intersection
of all type-$T$ congruences $\omega$ with $\cover{\omega}\supseteq
(B_{1}'')^{2}$. Every such $\omega$ has $\cover{\omega}\supseteq B_{1}^{2}$,
so that family is a \emph{subfamily} of $\{\omega_{1},\dots,\omega_{s}\}$ and
$\Omega''\supseteq\Omega$; hence $(B_{1}'\circ\Omega)\cap D=\varnothing$ as
well. Finally, by \eqref{eq:mal-incl} and $C_{1}\subseteq B_{1}\subseteq
B_{1}'$,
\[
  (C_{1}\circ\Omega)\cap B_{2}\;\subseteq\;
  (B_{1}'\circ\Omega)\cap\bigcap_{i\in[t]}K_{i}\cap B_{2}\;=\;
  (B_{1}'\circ\Omega)\cap D\;=\;\varnothing . \qedhere
\]
\end{proof}

\begin{hazard}[Claim~\ref{cl:multiply-step} is not proved, and it is the whole
  difficulty]\label{haz:multiply-citation}
Everything above is a proof except the claim, and the claim carries two
separate obligations. Neither is discharged here, so
Lemma~\ref{lem:multiply-all-linear} is \emph{not} available; see
\S\ref{sec:status}.

\emph{(a) The step is not of type $\TBA$, $\TC$ or $\TS$.} For $\TS$ this is
Lemma~\ref{lem:intersection-good}(s): $B_{1}'$ then contains some
$G\subt{\TBA,\TC}B_{1}''$, and (s) applied to $B_{1}''$ and $D$ --- which meet,
by \eqref{eq:mal-step} --- gives $G\cap D\neq\varnothing$, contradicting
$B_{1}'\cap D=\varnothing$. For $\TBA$ and $\TC$ separately the same route is
\emph{not} available, because $\subt{\TBA,\TC}$ in (s) and in
Lemma~\ref{lem:strong-nonempty} is a conjunction
(Caveat~\ref{haz:ba-central-meets}) and a single absorbing step supplies only
one of the two types. Lemma~\ref{lem:intersect-all}(t) gives
$B_{1}'\cap D\dsubte{\TBA}B_{1}''\cap D$, whose dot is exactly the possibility
to be excluded.

\emph{(b) When $T\in\{\TL,\TPC\}$, a type-$\TD$ step is not enough.} The
inductive hypothesis needs $B_{1}'\subt{\MT T}B_{1}''$ for the \emph{same} $T$.
A $\TD$ step carries an irreducible congruence, which is linear or PC by
Caveat~\ref{haz:linear-def}, but nothing so far forces it to match $T$. This
obligation is invisible while the components are written as $\TD$, which is why
the typing in the first paragraph of the proof matters.

Obligation (b) is the more serious: it is not a suppressed step but a missing
hypothesis in the induction, and closing it plausibly means strengthening the
inductive statement to carry the types of the intervening steps rather than
patching the claim.
\end{hazard}

\begin{proof}[Proof of Lemma~\ref{lem:maximal-mult}]
Write $\bar X$ for $X/\sigma$ and $\bar{\mathbf A}=\mathbf A/\sigma$. Note
first that for any $X\subseteq A$, $(X\circ\sigma)\cap B_{2}=\varnothing$ says
exactly that no $\sigma$-class meets both $X$ and $B_{2}$, i.e.\ that
$\bar X\cap\bar B_{2}=\varnothing$. So the hypotheses read
$\bar C_{1}\cap\bar B_{2}=\varnothing$ and, by
\ref{hyp:maximal-mult}, $\bar P\cap\bar B_{2}=\varnothing$; and
$\bar B_{1}\cap\bar B_{2}\neq\varnothing$ because
$B_{1}\cap B_{2}\neq\varnothing$.

\emph{Step 1: images commute with intersection on $F$.} For $j\in F$ the factor
$\bar C_{1}^{j}$ lands in the type-$T$ alternative of
Lemma~\ref{lem:propagation}(ft), so the ``moreover'' clause of
Lemma~\ref{lem:factor-by-delta} gives
$\sigma\cap B_{1}^{2}\subseteq\sigma_{j}\cap B_{1}^{2}$. Hence
\[
  \bigcap_{j\in F}\bar C_{1}^{j}\;=\;\bar P .
\]
Indeed $\subseteq$ is the only direction needing an argument. Let
$[a]\in\bigcap_{j\in F}\bar C_{1}^{j}$; since each
$\bar C_{1}^{j}\subseteq\bar B_{1}$ we may pick $a\in B_{1}$. For each
$j\in F$ there is $a_{j}\in C_{1}^{j}$ with $(a,a_{j})\in\sigma$; as
$a,a_{j}\in B_{1}$ this puts $(a,a_{j})\in\sigma\cap B_{1}^{2}\subseteq
\sigma_{j}$, so $a$ lies in the same $\sigma_{j}$-block as $a_{j}$, and
$C_{1}^{j}=B_{1}\cap E_{j}$ for that block gives $a\in C_{1}^{j}$. So
$a\in P$.

\emph{This is what \S\ref{sec:legislation}, item \ref{leg:quotient} warns
about, and it is exactly where the warning does not bite}: among the type-$T$
factors, and only among them, the saturation identity holds and the image of
the intersection is the intersection of the images.

\emph{Step 2: a proper multi-type reduction modulo $\sigma$.} Each
$\bar C_{1}^{j}\subt[\bar A]{T}\bar B_{1}$ for $j\in F$, and $\bar P$ is
nonempty because $P\supseteq C_{1}\neq\varnothing$. So by Step 1
\[
  \bar P\subte[\bar A]{\MT T}\bar B_{1},
\]
and the containment is \emph{proper}, since
$\bar P\cap\bar B_{2}=\varnothing$ while
$\bar B_{1}\cap\bar B_{2}\neq\varnothing$. Also $\bar B_{1}\lll\bar A$ and
$\bar B_{2}\lll\bar A$ by Corollary~\ref{cor:prop-quotient}(f).

\emph{Step 3: multiply by all type-$T$ congruences.} Apply
Lemma~\ref{lem:multiply-all-linear} inside $\bar{\mathbf A}$ to
$\bar P\subt[\bar A]{\MT T}\bar B_{1}\lll\bar A$ and
$\bar B_{2}\lll\bar A$: it returns the congruences
$\delta_{1},\dots,\delta_{s}$ of type $T$ on $\bar{\mathbf A}$ with
$\cover{\delta_{i}}\supseteq\bar B_{1}^{2}$, and
$\bigl(\bar P\circ\bigcap_{i}\delta_{i}\bigr)\cap\bar B_{2}=\varnothing$.
Note the ambient here is still $\bar B_{1}$, which is what the conclusion's
condition $\cover{\omega_{i}}\supseteq B_{1}^{2}$ requires.

\emph{Step 4: back to $\mathbf A$.} Let $\omega_{i}$ be the preimage of
$\delta_{i}$ under $\mathbf A\twoheadrightarrow\bar{\mathbf A}$, so
$\omega_{i}\supseteq\sigma$. By Lemma~\ref{lem:irred-quotient} the
correspondence $\omega\mapsto\omega/\sigma$ preserves irreducibility,
linearity and PC-ness and carries covers to covers, so each $\omega_{i}$ is of
type $T$ with $\cover{\omega_{i}}\supseteq B_{1}^{2}$. Pulling Step 3 back,
$\bigl(P\circ\bigcap_{i}\omega_{i}\bigr)\cap B_{2}=\varnothing$, and
$C_{1}\subseteq P$ gives
$\bigl(C_{1}\circ\bigcap_{i}\omega_{i}\bigr)\cap B_{2}=\varnothing$.

Finally $\bigcap_{i}\omega_{i}\supseteq\sigma$, and it has the property that
$\sigma$ was chosen maximal for; so $\sigma=\bigcap_{i}\omega_{i}$.
\end{proof}

\begin{hazard}[What Hypothesis~\ref{hyp:maximal-mult} costs, and three failed
  routes]\label{haz:maximal-mult}
What the other hypotheses give is
$\bar C_{1}\cap\bar B_{2}=\varnothing$, and $C_{1}\subseteq P$, so
$\bar C_{1}\subseteq\bar P$: Hypothesis~\ref{hyp:maximal-mult} asks for
emptiness of the \emph{larger} set, and is therefore strictly stronger. Three
routes to it have been tried and none works.

\emph{Drop the non-type-$T$ factors one at a time.} If $j_{0}\notin F$ and
$C_{1}'=\bigcap_{j\neq j_{0}}C_{1}^{j}$, one would like
$\bar C_{1}'\cap\bar B_{2}=\varnothing$ and then induct. But recovering that
from $\bar C_{1}\cap\bar B_{2}=\varnothing$ needs
$\bar C_{1}=\bar C_{1}'\cap\bar C_{1}^{j_{0}}$, which is the
image-versus-intersection identity again --- available for type-$T$ factors by
Step~1 of the proof above and \emph{not} for the others, since the saturation
clause of Lemma~\ref{lem:factor-by-delta} is exactly what the $\subt{\TS}$
alternative withholds.

\emph{Pass to a minimal subfamily.} Choosing $F'$ minimal with
$\bigcap_{j\in F'}\bar C_{1}^{j}\cap\bar B_{2}=\varnothing$ would make
minimality supply the leave-one-out hypothesis of
Theorem~\ref{thm:stable-intersection}, whose four conclusions have no room for
$\TS$ --- so no member of $F'$ would be $\TS$-typed. But no such $F'$ need
exist: the whole family gives
$\bigl(\bigcap_{j}C_{1}^{j}\bigr)/\sigma\cap\bar B_{2}=\varnothing$, whereas
$\bigl(\bigcap_{j}C_{1}^{j}\bigr)/\sigma\subseteq\bigcap_{j}\bar C_{1}^{j}$
with the inclusion generally strict (\S\ref{sec:legislation}, item
\ref{leg:quotient}). This is the route an earlier draft of this document took,
and it is invalid.

\emph{Descend into a strong subuniverse.} The $\subt{\TS}$ factors do at least
never help: a factor with $\bar C_{1}^{j}\subt{\TS}\bar B_{1}$ contains a
subuniverse $\bar D$ of $\bar B_{1}$ that is simultaneously binary absorbing
and central, and Lemma~\ref{lem:strong-nonempty} makes $\bar D$ meet
$\bar B_{1}\cap\bar B_{2}$, so $\bar C_{1}^{j}\cap\bar B_{2}\neq\varnothing$.
But the configuration reproduces itself on $\bar D$ --- $\bar D\lll\bar A$,
$\bar C_{1}\cap\bar D\lll\bar D$, $\bar B_{2}\cap\bar D\lll\bar D$, still
disjoint, with $\bar D\cap\bar B_{2}\neq\varnothing$ --- so no iteration of
that fact closes anything, and the descent shrinks the ambient while the
conclusion's congruence family is indexed by the ambient.

So the hypothesis stays, and the debt is recorded in \S\ref{sec:status}.

A fourth route is open, and unlike the three above it has not failed --- it
has not been attempted. Lemma~\ref{lem:maximal-mult} has exactly one consumer,
the typing step of Lemma~\ref{lem:parallelogram-crucial}, and the older proof
of the dichotomy establishes the parallelogram property of a crucial
constraint \emph{without any statement of this shape}: by recursion on a
measure over reductions, restricting to the minimal linear reduction
containing a solution of the weakened instance, with a fit-in lemma for
one-of-four chains through a subdirect relation doing the intersection work.
No maximal congruence appears. If Lemma~\ref{lem:parallelogram-crucial} is
proved along those lines, its typing conclusion recovered separately ---
irreducibility from Lemma~\ref{lem:crucial-irreducible}, the type from the
bridge machinery --- then the sole call site of
Lemma~\ref{lem:maximal-mult} disappears, and
Hypothesis~\ref{hyp:maximal-mult} with it. The transport cost is not small:
the fit-in lemma has to be restated in this document's vocabulary, and the
recursion needs the reduction machinery of \S\ref{sec:main-statements}. But it
converts an open problem into writing, and it is the first route with that
property.
\end{hazard}
